\documentclass[11pt]{article}

\usepackage[utf8]{inputenc}
\usepackage[T1]{fontenc}
\usepackage{lmodern}
\usepackage[margin=1in]{geometry}
\usepackage{amsmath,amssymb,amsthm,mathtools}
\usepackage{booktabs}
\usepackage{siunitx}
\usepackage{xcolor}
\usepackage{algorithm}
\usepackage{algpseudocode}
\usepackage{pgfplots}
\pgfplotsset{compat=1.18}
\usepgfplotslibrary{groupplots}
\usepackage[colorlinks=true,linkcolor=blue!50!black,citecolor=blue!50!black,urlcolor=blue!50!black]{hyperref}
\usepackage{cleveref}

% ------------------------------------------------------------------ theorems
\newtheorem{proposition}{Proposition}[section]
\newtheorem{lemma}[proposition]{Lemma}
\theoremstyle{definition}
\newtheorem{definition}[proposition]{Definition}
\theoremstyle{remark}
\newtheorem{remark}[proposition]{Remark}

% ------------------------------------------------------------------ notation
\DeclareMathOperator{\Diag}{Diag}
\DeclareMathOperator{\diag}{diag}
\DeclareMathOperator{\rank}{rank}
\DeclareMathOperator{\range}{range}
\newcommand{\Sp}{\mathcal{S}_+^n}
\newcommand{\Proj}{\Pi_{\mathcal{S}_+^n}}
\newcommand{\Xs}{X^\star}
\newcommand{\ys}{y^\star}
\newcommand{\fro}[1]{\lVert #1\rVert_F}
\newcommand{\nrm}[1]{\lVert #1\rVert}
\newcommand{\solver}[1]{\textsf{#1}}
\newcommand{\EVD}{\textsc{evd}}
\newcommand{\EVDs}{\textsc{evd}s}
\newcommand{\sep}{\mathrm{sep}}
\newcommand{\mubh}{\mu_{\mathrm{BH}}}

\title{Exact test problems and a work--accuracy protocol\\ for benchmarking nearest correlation matrix solvers}
\author{Oyeyinka E. Ibrahim \\
  Department of Industrial and Systems Engineering \\
  Northern Illinois University \\
  \texttt{ibrahimoyeyinka@gmail.com}}
\date{Draft of \today}

\begin{document}
\maketitle

\begin{abstract}
Comparisons of nearest correlation matrix (NCM) solvers are usually run on
collections of matrices whose solutions are unknown, so the accuracy each
method actually delivers is not measured. We specialize the known technique of
generating test problems from prescribed optimality conditions to the NCM
problem, obtaining instances whose primal and dual solutions are known in
closed form and whose solution rank, near-zero multiplicity and eigenvalue
separation are set independently. Measuring forward error against those exact
solutions changes what a comparison reports. The semismooth Newton method with
the Borsdorf--Higham line search is cheapest under all 11 conventions we
examine, on every test set, and its cost is flat in dimension: five
eigendecompositions at $n=100$ and again at $n=2105$. The ordering below it is
a property of the test set rather than of the methods. It changes with the
accuracy target, with how far the input is from a correlation matrix, and above
all with the solution rank, whose crossover we bracket between $r=5$ and $r=20$
at $n=100$ and between $r=50$ and $r=75$ at $n=500$, rising further but
unresolved on an equity panel at $n=2105$. Two further results are practical.
The dimension-scaled stopping rule $\nrm{\nabla\theta}_2\le10^{-7}n$ permits
roughly 340 times less accuracy at $n=500$ than at $n=100$, so a larger problem
can appear easier. And a known finite-precision failure of objective-value
Armijo tests recurs in a second published method, where it does not affect the
published results but stalls the method on real and equity-derived matrices up
to $n=3250$; a safeguard adapted from the NCM literature removes it in every
case.
\end{abstract}
% Word count of abstract: approx. 215.
% Corrected 2026-09-11: the per-cell count is 7 orderings (not 9) under the
% corrected reach-and-hold rule; the rank attribution replaces it in the abstract.

\tableofcontents

\section{Introduction}\label{sec:intro}

Given a symmetric matrix $G\in\mathbb{R}^{n\times n}$, the nearest correlation
matrix problem asks for
\begin{equation}\label{eq:ncm}
  \min_{X}\ \tfrac12\fro{X-G}^2 \quad\text{subject to}\quad X\succeq 0,\ \diag(X)=e .
\end{equation}
The problem has a unique solution. Since Higham's alternating-projections
method~\cite{higham2002nearest}, the literature has produced dual formulations and
first-order schemes~\cite{malick2004dual,boyd2005least}, a quadratically convergent
semismooth Newton method~\cite{qisun2006quadratically}, its preconditioned and globalized
refinement~\cite{borsdorf2010preconditioned}, accelerated projection
methods~\cite{higham2016anderson}, and, most recently, accelerated gradient
methods with quasi-Newton preconditioning~\cite{huynh2025accelerated}. Most of these
papers include numerical comparisons, and those comparisons are how a reader
decides which method to use.

This paper is about how such comparisons are made. We ask a narrow question:
\emph{when NCM solvers are compared, which methodological choices change the
conclusion, and which only change the numbers?} We identify four candidate
confounds:
\begin{enumerate}
  \item \emph{Native stopping rules.} Each method stops when its own criterion
    is met, so the solvers stop at different accuracies.
  \item \emph{Approximate references.} Forward error is measured against the
    output of another solver instead of against the true solution.
  \item \emph{Rejected-iterate accounting.} Line-search trials that are
    rejected can be credited with reaching a target that the method has not
    actually reached.
  \item \emph{Hidden primitive work.} `Iterations' count different
    operations in different methods, and some per-iteration work is not
    recorded at all.
\end{enumerate}
A careful benchmark must control each of these. Without an exact solution,
however, the second confound cannot be removed, and without trajectories the
first and third cannot be separated.

\paragraph{Contributions.}
In brief, we give an exact-solution NCM family and a matched-accuracy
protocol, and use them to show that conclusions drawn from NCM solver
comparisons do not transfer across benchmark regimes: the rank at which an
ordering established at $n=100$ reverses moves higher at $n=500$, the order
also changes with the degree of invalidity on matrices built from equity
returns, and four matrices from the literature are consistent with this
dependence. We also show that a known
finite-precision line-search defect recurs in a second published method and
occurs on real matrices.
Each contribution builds on established practice (\cref{sec:related}).
\begin{enumerate}
  \item \emph{An exact-solution NCM family} (\cref{sec:family}). We specialize
    the standard construction of test problems from prescribed optimality
    conditions~\cite{lenard1984randomly,calamai1993new} to the NCM problem.
    For semidefinite programs, the same idea is due to Wei and
    Wolkowicz~\cite{wei2010generating}. The parts specific to NCM are the
    projection identity $\Proj(G+\Diag\ys)=\Xs$, independent control of the
    rank, the multiplicity of the smallest eigenvalue of $C(\ys)$ and its
    distance from zero (the strict-complementarity defect equals the
    multiplicity when that distance is zero and vanishes otherwise, so 45 of
    the 270 instances at each dimension are not strictly complementary), a paired-seed
    design, sufficient conditions for $|G_{ij}|\le1$ together with a
    fixed-rank saturation result for $\max_{i\ne j}|\Xs_{ij}|$, and released
    instance sets at $n=100$ and $n=500$.
  \item \emph{A matched-accuracy protocol for NCM} (\cref{sec:protocol}). We
    adapt data profiles~\cite{more2009benchmarking}, fixed-target
    evaluation~\cite{hansen2021coco} and work--precision
    diagrams~\cite{hairer1996solving} to spectral solvers. The unit of cost is
    one \EVD, which we check against wall-clock time. Because the reference is
    exact, two rules become enforceable: rejected trials receive no accuracy
    credit, and accuracy is matched in forward error. Primitive work is
    reported separately.
  \item \emph{An invariant first place above a regime-dependent ordering}
    (\cref{sec:ranking}). Six verified variants of five methods were run on
    270 instances at $n=100$ and 270 at $n=500$, with 135 further instances at
    higher rank, 125 built from two equity panels at $n=550$ and $n=2105$, and
    four matrices from the literature. The cheapest method is the same under
    all 11 conventions examined, and its cost is flat in dimension: 5
    eigendecompositions at $n=100$ and at $n=2105$. Second at $n=100$ is
    Anderson-accelerated alternating projections~\cite{higham2016anderson},
    under 10 of the 11 conventions and tied for second under the eleventh; it
    keeps that place on the exact-solution families but not on perturbed or
    literature matrices.
    Below the first place the ordering changes with the solution rank and
    with how far the input is from a correlation matrix, and, for SBB-Dual
    against AGD-SDAJ-BH, with the accuracy target. The rank dependence holds
    between published methods alone:
    Dykstra-APM and AGD-SDAJ-BH exchange places between $r=20$ and $r=50$ at
    $n=100$, and the exchange has not occurred by $r=150$ at $n=500$. At
    $n=2105$ the order of Anderson-APM and AGD-SDAJ-BH is not monotone in
    rank: Anderson-APM is cheaper at $r=50$ and $r=400$, AGD-SDAJ-BH at $r=100$
    and $r=200$. The same rank pattern appears between SBB-Dual and
    AGD-SDAJ-BH, whose exchange is bracketed between $r=5$ and $r=20$ at
    $n=100$ and between $r=50$ and $r=75$ at $n=500$. Only a family that varies
    the spectral structure at the solution in controlled steps makes such an
    association directly observable.
  \item \emph{Replication of a known line-search defect}
    (\cref{sec:traps}). The failure of objective-value sufficient-decrease
    tests near a minimizer is well known~\cite{hager2005new}. For the NCM dual,
    Borsdorf analysed it~\cite[\S4.7.1]{borsdorf2007newton}, and this led to
    the modified step selection of Borsdorf and
    Higham~\cite{borsdorf2010preconditioned}. The new content is limited to
    two points. First, the defect recurs in a separately published
    first-order method~\cite{huynh2025accelerated}. Second, we quantify it on
    the controlled family: 56/270 stalls for semismooth Newton and 18/270 for
    AGD-SDAJ, both removed by the Borsdorf--Higham correction. We also
    document a second, simpler trap (a redundant Jacobian factorization).
  \item \emph{An NCM quantification of the stopping-rule confound}
    (\cref{sec:stopping}). Dolan, Mor\'e and Munson~\cite{dolan2006optimality}
    warned that native stopping tests confound comparisons. For NCM, the
    widely used rule $\nrm{\nabla\theta}_2\le 10^{-7}n$ leaves the achieved
    forward error spread over a factor of 383 across solvers at $n=100$, and
    it degrades accuracy with dimension unequally across methods.
\end{enumerate}

\paragraph{What we do not claim.}
We do not claim that any method is best in general. The ranking results are
for a synthetic family at $n=100$ and $n=500$, with instances built from
equity returns at $n=550$ and $n=2105$ (\cref{sec:thesis,sec:us2105}) and four
matrices from the literature (\cref{sec:real}) as external checks, all on one
machine. SBB-Dual is a new method from the author's companion manuscript in
preparation~\cite{ibrahim2026sbb}. It is specified in full in
\cref{app:algorithms}, receives the same exposition as the others, and its
convergence theory is not used here. Three findings do not depend on it:
without SBB-Dual, Newton-SIN-BH is still first, Anderson-APM is still
second at $n=100$, and the rank dependence is still present, between
Dykstra-APM and AGD-SDAJ-BH. Two findings do use it: the dependence of the
order on the accuracy target, and the bracketed positions of the rank
crossover at $n=100$ and $n=500$ (\cref{sec:cells,sec:highrank}).

\subsection{Related work}\label{sec:related}

\paragraph{Benchmarking methodology.}
Performance profiles~\cite{dolan2002benchmarking} compare solvers by the ratio
of their cost to the best cost on each problem. Cost is measured at each
solver's own termination, so the solver's stopping test implicitly sets the
accuracy. Dolan, Mor\'e and Munson~\cite{dolan2006optimality} showed that
differing stopping criteria bias such profiles, and proposed rerunning under a
common optimality measure. Data profiles~\cite{more2009benchmarking} count
cost in a problem-independent unit and show how rankings change with the
required accuracy. COCO~\cite{hansen2021coco} uses a ladder of fixed targets
against known optima. Work--precision diagrams, standard in the comparison of
ODE solvers~\cite{hairer1996solving}, plot error against a reference solution
versus work. Gould and Scott~\cite{gould2016note} show that performance-profile
rankings below the leading solver can change when solvers are added or
removed. Beiranvand, Hare and Lucet~\cite{beiranvand2017best} survey best
practice, including known-solution test problems and reproducibility. Our
protocol follows these principles; our contribution is to make their
exact-reference versions possible for the NCM problem. The second/third-place
reversals of \cref{sec:ranking} arise with a \emph{fixed} solver set, from the
accuracy target and the instance regime, so the mechanism differs from the one
studied by Gould and Scott.

\paragraph{Test problems with known solutions.}
Generating test problems from prescribed optimality conditions is standard
for unconstrained optimization~\cite{more1981testing} and for quadratic
programming~\cite{lenard1984randomly,calamai1993new}. Wei and
Wolkowicz~\cite{wei2010generating} generate semidefinite programs from a
chosen optimal primal--dual pair with prescribed complementarity nullity, and
show that nullity correlates with solver difficulty. Our family applies the
same idea to the NCM problem. Their generator does not apply directly: it
builds a linear-objective SDP and chooses the constraint operator as part of
the construction, whereas NCM fixes the constraint operator as $\diag(X)=e$
and leaves only $G$ free. As far as we have found, no NCM study uses
instances with a known exact solution. The test sets
of~\cite{higham2002nearest,qisun2006quadratically,borsdorf2007newton,borsdorf2010preconditioned,higham2016anderson,armijo2025semismooth}
are random or data-derived, as are the Anymatrix \texttt{CORRINV}
matrices~\cite{higham2022anymatrix} used in~\cite{huynh2025accelerated}, and
in all of them the reference solutions are computed approximately.

\paragraph{Finite-precision line searches.}
An objective-value sufficient-decrease test cannot be evaluated accurately near
a minimizer, because the difference of function values suffers cancellation.
Hager and Zhang~\cite{hager2005new} introduced approximate Wolfe conditions
for this reason. For the NCM dual, Borsdorf~\cite[\S4.7.1]{borsdorf2007newton}
derived the stagnation threshold, $\nrm{\nabla\theta}$ of the order of the
square root of the rounding error in $\theta$ (about $10^{-8}$), noted that
the problem is general to Armijo backtracking, and proposed the alternatives
that became the modified step selection of~\cite{borsdorf2010preconditioned}.
\Cref{sec:traps} replicates this known phenomenon; it does not rediscover it.

\subsection{Outline}
\Cref{sec:family} constructs the instance family. \Cref{sec:protocol} states
the protocol, \cref{sec:traps} describes the two implementation traps,
\cref{sec:ranking} reports the ranking study, and \cref{sec:stopping}
examines the dimension-scaled stopping rule. \Cref{sec:repro} covers
reproducibility and \cref{sec:discussion} discusses limitations. The
appendices contain the proofs for the instance family, the verification
records for our reimplementations of AGD-SDAJ and Anderson-APM, the artifact
manifest, and pseudocode for every solver.

\section{The KKT-controlled instance family}\label{sec:family}

\subsection{Dual formulation}
Following~\cite{malick2004dual,qisun2006quadratically}, the Lagrangian dual of~\eqref{eq:ncm}
is the unconstrained problem
\begin{equation}\label{eq:dual}
  \min_{y\in\mathbb{R}^n}\ \theta(y) \coloneqq \tfrac12\fro{\Proj(G+\Diag(y))}^2 - e^{\mathsf T}y ,
  \qquad
  \nabla\theta(y) = \diag\bigl(\Proj(G+\Diag(y))\bigr) - e .
\end{equation}
Here $\Proj$ is the projection onto the positive semidefinite cone. The
gradient is globally Lipschitz with constant 1. If $\ys$ satisfies
$\nabla\theta(\ys)=0$, then $\Xs=\Proj(G+\Diag(\ys))$ is the unique solution
of~\eqref{eq:ncm}. Each evaluation of $\theta$ or $\nabla\theta$ requires one
symmetric eigendecomposition (\EVD) of $C(y)\coloneqq G+\Diag(y)$.

\subsection{Construction and exactness}\label{sec:construction}
The idea is to choose the solution first and then build a $G$ for which that
solution satisfies the KKT conditions exactly. This is the standard technique
of generating instances from prescribed optimality
conditions~\cite{lenard1984randomly,calamai1993new}, and in particular Wei and
Wolkowicz's generator of semidefinite programs with prescribed complementarity
nullity~\cite{wei2010generating}, specialized here to the NCM problem. It also
continues the long tradition of test problems with closed-form solutions
surveyed in Higham's gallery of test matrices~\cite[Ch.~26]{higham2002accuracy}. What is
specific to NCM is the exact projection identity of \cref{prop:exact}, the
control of the rank, multiplicity and separation in \cref{sec:spectrum}, and
the entrywise conditions of \cref{sec:feasibility}.
Their construction does not cover the present problem: it generates the data
of a \emph{linear}-objective SDP and is free to choose the constraint operator
itself (their Algorithm~2.3, Steps 3--4), whereas the NCM problem fixes both
the quadratic objective and the operator $\mathcal{A}=\diag$ with $b=e$, leaving
$G$ as the only free data. The shared element is the idea of prescribing an
optimal primal--dual pair and controlling the dimension of its common
nullspace, their complementarity nullity, which corresponds to $|\beta|$ here.

\begin{definition}[KKT-controlled instance]\label{def:family}
Let $n>r\ge1$. Draw $U\in\mathbb{R}^{n\times r}$ with independent
standard-normal entries and rescale every row to unit Euclidean norm. Set
$\Xs=UU^{\mathsf T}$, so that $\diag(\Xs)=e$, $\Xs\succeq0$ and $\rank\Xs=r$
almost surely. Let $Q_0\in\mathbb{R}^{n\times(n-r)}$ have orthonormal columns
spanning $\ker\Xs$. For $\mu\in\mathbb{R}^{n-r}_{\ge0}$ define
\begin{equation}\label{eq:construction}
  W = Q_0\Diag(\mu)Q_0^{\mathsf T},\qquad
  G = \Xs - W + \Diag(\diag W),\qquad
  \ys = -\diag(W).
\end{equation}
\end{definition}

\begin{proposition}[Exactness]\label{prop:exact}
For every $\mu\ge0$, $\Proj(G+\Diag(\ys))=\Xs$ and $\nabla\theta(\ys)=0$.
Hence $\Xs$ is the unique nearest correlation matrix to $G$.
\end{proposition}
\begin{proof}
By construction, $G+\Diag(\ys)=\Xs-W$. Both $\Xs$ and $W$ are symmetric
positive semidefinite, and $\Xs W=\Xs Q_0\Diag(\mu)Q_0^{\mathsf T}=0$ because
the columns of $Q_0$ lie in $\ker\Xs$. The ranges of $\Xs$ and $W$ are
therefore orthogonal. They need not be complementary: if some $\mu_j=0$, then
$\range W$ is a proper subspace of $\ker\Xs$. Consequently $\Xs-W$ has the
spectral decomposition obtained by combining those of $\Xs$ and $-W$ on
orthogonal subspaces. Its positive part is $\Xs$, and its nonpositive part is
$-W$. Thus $\Proj(\Xs-W)=\Xs$, and
$\nabla\theta(\ys)=\diag(\Xs)-e=0$. Since $\theta$ is convex and
differentiable, $\ys$ minimizes~\eqref{eq:dual}. Uniqueness of the primal
solution is classical~\cite{higham2002nearest,qisun2006quadratically}.
\end{proof}

In floating point, the generator recomputes the projection rather than relying
on the identity. Over the 270 released instances at $n=100$, the KKT residual
is at most $1.7\times10^{-14}$ and the exactness error
$\fro{\Proj(C(\ys))-\Xs}$ at most $8.4\times10^{-14}$.

\subsection{Prescribed spectral structure}\label{sec:spectrum}
Because the ranges are orthogonal, the spectrum of $C(\ys)=\Xs-W$ is the
multiset union
\begin{equation}
  \operatorname{spec}C(\ys) = \{\lambda_1(\Xs),\dots,\lambda_r(\Xs)\}\ \cup\ \{-\mu_1,\dots,-\mu_{n-r}\}.
\end{equation}
In the notation of~\cite{qisun2006quadratically}, let $\alpha,\beta,\gamma$ index the
positive, zero and negative eigenvalues of $C(\ys)$. Then
\begin{itemize}
  \item $|\alpha|=r$, so the rank of the solution is prescribed;
  \item $|\beta|=\#\{j:\mu_j=0\}$ exactly, so the multiplicity of zero
    eigenvalues at the solution is prescribed;
  \item an entry $\mu_j=\delta>0$ produces an eigenvalue $-\delta$, which moves
    toward the nonsmooth boundary of $\Proj$ as $\delta\to0$.
\end{itemize}
With the layout used below ($m$ entries of $\mu$ equal to $\delta$, the rest
equal to $\mu_{\mathrm{bulk}}>0$), the multiplicity and the separation are
therefore not varied independently in their effect on strict
complementarity: $|\beta|=m$ when $\delta=0$ and $|\beta|=0$ otherwise.

\paragraph{Terminology.}
Qi and Sun~\cite{qisun2006quadratically} showed that constraint
nondegeneracy always holds for the NCM problem, so the family does not, and
could not, vary constraint degeneracy. What it varies is \emph{failure of
strict complementarity}: zero eigenvalues of $C(\ys)$ ($|\beta|>0$), and
near-zero eigenvalues (small $\delta$) that approach it. Throughout this
paper, a \emph{degeneracy cell} means one $(r,m,\delta)$ combination in this
sense, and the \emph{strict-complementarity defect} means $|\beta|$.
The generator uses
$\mu=(\mu_{\mathrm{bulk}},\dots,\mu_{\mathrm{bulk}},\delta,\dots,\delta)$ with
$m$ copies of $\delta$ and $\mu_{\mathrm{bulk}}=0.1$. The separation parameter
is $\sep=\delta/\mu_{\mathrm{bulk}}$ (we avoid $\rho$, which denotes the
line-search backtracking factor in \cref{sec:traps,app:algorithms}). The released grid is
\[
  r\in\{5,20,50\},\quad m\in\{1,5,20\},\quad
  \delta\in\{10^{-2},10^{-4},10^{-6},10^{-8},10^{-10},0\},
\]
with five paired seeds per rank. This gives $3\times3\times6=54$ cells and
270 instances at $n=100$. In plots, $\sep$ is shown on a $\log_{10}$ axis and
the case $\sep=0$ is shown separately.

\subsection{Correlation-like inputs: sufficient conditions for $|G_{ij}|\le1$}\label{sec:feasibility}
\Cref{prop:exact} places no restriction on the entries of $G$. For the
instances to resemble invalid correlation matrices, which is the setting in
which NCM problems arise, the generator additionally requires
$\diag(G)=e$ (true by construction), $\lambda_{\min}(G)<0$, and
$|G_{ij}|\le1$. The first two are checked directly. The results below explain
when the third can be expected to hold. Let $P_r$ be the orthogonal projector
onto $\range U$, so that $Q_0Q_0^{\mathsf T}=I-P_r$, and define
\[
  a \coloneqq \max_{i\ne j}|\Xs_{ij}|,\qquad b\coloneqq \max_{i\ne j}|(P_r)_{ij}|.
\]

\begin{proposition}[Deterministic bound, uniform $\mu$]\label{prop:det}
If $\mu=c\,e$ with $c>0$, then $|G_{ij}|\le a+cb$ for all $i\neq j$. In
particular, if $a<1$, then $0<c\le(1-a)/b$ is sufficient for $|G_{ij}|\le1$.
\end{proposition}

\begin{proposition}[Fixed-rank saturation]\label{prop:sat}
Fix $r\ge2$ and let the rows $u_i$ of $U$ be independent and uniform on
$S^{r-1}$. Then $a\to1$ as $n\to\infty$, along every realization.
\end{proposition}

\Cref{prop:sat} says only that $a$ saturates. It does \emph{not} imply that a
fixed $c>0$ eventually violates \cref{prop:det}: at fixed $r$ the quantity $b$
is observed to decrease with $n$ as well, but we do not prove
this, so whether $a+cb$ exceeds 1 depends on the relative
rates of $1-a$ and $b$, which we do not establish. We therefore make no
asymptotic feasibility claim, and the generator relies on entrywise
screening.

\begin{proposition}[High-dimensional regime]\label{prop:hd}
Let the rows of $U$ be independent and uniform on $S^{r-1}$ and let
$\eta\in(0,1)$ be a failure probability (distinct from the near-zero value
$\delta$ in $\mu$). Then with probability at least
$1-\eta$,
\[
  a\;\le\;\sqrt{\frac{2}{r}\Bigl(2\log n+\log\frac{2}{\eta}\Bigr)} .
\]
In particular $a\to0$ in probability whenever $r/\log n\to\infty$.
\end{proposition}

The bound follows from the explicit cap estimate
of~\cite[Lem.~2.2]{ball1997elementary} (subgaussianity in the sense
of~\cite[Thm.~3.4.5]{vershynin2026high} gives the same tail without a
specific constant) and a union bound; see
\cref{app:family}. We state it only for $a$, the quantity that governs
\cref{prop:det} in the regime of interest. A companion bound on $b$ would
require a model for $\range U$ that we do not need here, since the released
family lies in the fixed-rank regime of \cref{prop:sat}, where
\cref{prop:hd} does not apply.

Proofs are given in \cref{app:family}. These results give sufficient
conditions only. Failure of the bound in \cref{prop:det} does not show that
some $|G_{ij}|$ exceeds 1, and for that reason the generator checks every
entry directly instead of relying on the bound. \Cref{prop:det} assumes
uniform $\mu$, whereas the released instances use a mixed $\mu$ (bulk value
$0.1$ together with $m$ copies of $\delta$).
The extension to mixed $\mu$ is left open: the extra off-diagonal term it
introduces depends on the choice of orthonormal basis for $\ker\Xs$ and is not
determined by $(n,r,m,\delta)$ alone (\cref{app:family}). The released
mixed-$\mu$ instances are therefore screened entrywise rather than certified
by a bound.
None of these results says anything about indefiniteness of $G$, which is
checked separately.

\paragraph{Which regime the released family is in.}
\Cref{tab:ab} reports empirical values of $(a,b)$ at $n=100$ over 30 seeds per
rank. The family is in the regime of \cref{prop:sat}, not \cref{prop:hd}. The
rank is fixed and small relative to $n$, $a$ is close to 1 at $r=5$, and the
bound on $c$ implied by \cref{prop:det} shrinks sharply as the rank decreases.

\begin{table}[t]
\centering
\caption{Empirical $(a,b)$ at $n=100$, 30 seeds per rank. The last column is
the sufficient bound $(1-a)/b$ of \cref{prop:det} evaluated at the mean values
of $a$ and $b$.}
\label{tab:ab}
\small
\begin{tabular}{rrrrrr}
\toprule
$r$ & mean $a$ & max $a$ & mean $b$ & max $b$ & $(1-\bar a)/\bar b$\\
\midrule
5  & 0.9905 & 0.9981 & 0.0590 & 0.0708 & 0.162\\
20 & 0.7495 & 0.8340 & 0.1483 & 0.1871 & 1.690\\
50 & 0.5081 & 0.5719 & 0.1856 & 0.2083 & 2.650\\
\bottomrule
\end{tabular}
\end{table}

The table is consistent with the screening outcomes. At
$\mu_{\mathrm{bulk}}=0.1$, the sufficient bound at the mean values at $r=5$
($0.162$) exceeds $0.1$, but at the worst seeds (max $a=0.9981$) the bound
itself falls below $0.1$. In the
paired design (\cref{sec:paired}), 5 of the 7 candidate seeds were accepted at
$r=5$; the 2 rejections were for $|G_{ij}|>1$. At $r=20$ and $r=50$, 5 of 5
were accepted. With $\mu_{\mathrm{bulk}}=1$, the sufficient bound at the mean
values still holds at $r=20$ and $r=50$ ($1.690$ and $2.650$ exceed 1) but
fails at $r=5$, and in an earlier design every $r=5$ candidate was rejected.
Since one bulk level is used at every rank, this is why the generator uses
$\mu_{\mathrm{bulk}}=0.1$.

\subsection{Paired-seed design and screening}\label{sec:paired}
Each candidate seed is screened once per rank against its entire
$(m,\delta)$ block. A seed is accepted only if all $3\times6$ instances built
from it pass the checks ($\diag(G)=e$, $|G_{ij}|\le1$, $\lambda_{\min}(G)<0$,
and a recomputed KKT residual and exactness error each below $10^{-9}$). The
accepted seeds are then reused across every $(m,\delta)$ cell. Varying the
separation therefore changes only $\mu$, with $U$ held fixed, so comparisons
along $\delta$ are within-seed. Rejections are counted and reported, not
silently discarded.

Because of the screening, the $r=5$ sample is conditioned differently from the
$r=20$ and $r=50$ samples: at $r=5$ it excludes the seeds whose instances have
an off-diagonal entry above 1 (at $n=100$, both rejections occurred at $m=20$,
$\delta=10^{-2}$). These are not simply the seeds with the largest $a$, since
the rejection depends on the basis-dependent kernel term as well.
Comparisons across ranks should be read with this in mind.

\paragraph{Extent of the family.}
Every synthetic instance is close to feasible: $\lambda_{\min}(G)$ lies between
$-0.075$ and $-0.017$ at $n=100$ and between $-0.047$ and $-0.011$ at $n=500$.
A larger bulk level would give more strongly invalid inputs, but it fails the
entrywise screen at $r=5$, and the family uses a single bulk level at every
rank; the degree of invalidity is therefore not a controlled factor here.
Findings that depend on it (\cref{sec:thesis,sec:us2105}) rest on perturbed
matrices whose reference is computed, not exact.

\paragraph{Regeneration.}
The seed of candidate $k$ at rank $r$ and dimension $n$ is $100000n+1000r+k$.
The multiplicity is limited by the kernel dimension, $m\le n-r$, and the $m$
entries equal to $\delta$ are attached to the last $m$ columns of $Q_0$, the
kernel block of the eigenvector basis of $\Xs$ returned by the symmetric
eigensolver, which is the point at which a
different LAPACK build can yield a different member of the family
(\cref{app:artifact}).

\subsection{Release}
We release the 270 paired instances ($G$, $\Xs$, $\ys$, SHA-256 hashes of the
upper triangle of $G$, and per-instance screening diagnostics), the generator,
and the seed-screening log. \Cref{app:artifact} lists the files.
We also release a paired set at $n=500$, generated on the same grid. All 270
instances passed screening; seed acceptance was 5 of 7 candidates at $r=5$
(both rejections for $|G_{ij}|>1$) and 5 of 5 at $r=20$ and $r=50$, with
$\max_{i\ne j}|G_{ij}|=0.99987$ among the accepted $r=5$ instances. The
ranking study of \cref{sec:n500} uses all five paired seeds per rank.

\section{Protocol}\label{sec:protocol}

This section sets out the measurement protocol. Each rule addresses a
confound that we observed on this family, and each is stated together with the
evidence that motivates it.

\subsection{Native stopping rules as a confound}\label{sec:native}
Suppose every solver is run to its own exit criterion and the resulting costs
are compared. The comparison then credits a method for stopping early as much
as for converging quickly. This happens even when all methods share one
nominal rule. Under the dimension-scaled rule $\nrm{\nabla\theta}_2\le10^{-7}n$
used by~\cite{huynh2025accelerated}, which at $n=100$ is $10^{-5}$, the six
solvers of \cref{sec:ranking} and the Newton-SIN control stop at median
forward errors between
$5.04\times10^{-8}$ and $1.93\times10^{-5}$, a factor of 383
(\cref{sec:stopping}, \cref{fig:native}). A cost comparison taken at these
exits compares different accuracies.

\subsection{Work--accuracy trajectories against the exact solution}\label{sec:traj}
We take one \EVD\ as the unit of cost. After \emph{every} \EVD, whether it
produces an iterate, a line-search trial or a Jacobian factorization, we record
the pair
\[
  \bigl(\text{cumulative \EVDs},\ \fro{X_k-\Xs}\bigr),
\]
together with $\nrm{\nabla\theta}_2$ and a flag indicating whether the
evaluated point was accepted. Here $\Xs$ is the exact solution from
\cref{prop:exact}, not the output of another solver. Measuring forward error
rather than a residual follows the standard distinction
of~\cite[Ch.~1--2]{higham2002accuracy}. All runs use a common,
tight tolerance of $10^{-11}$, and each solver has its own cap (\cref{sec:ranking}). These trajectories are work--precision diagrams in the sense
of~\cite{hairer1996solving}, and reading costs off them at a ladder of targets
follows data profiles~\cite{more2009benchmarking} and fixed-target
evaluation~\cite{hansen2021coco}. The difference is that the reference is the
exact solution, not the best value found by any solver. For any target
$\varepsilon$, the cost can then be read off the trajectory after the run. We
define the \emph{cost to $\varepsilon$} as the cumulative \EVD\ count at the
first accepted iterate from which $\fro{X_k-\Xs}\le\varepsilon$ holds for
every later accepted iterate, that is, at the start of the final sub-$\varepsilon$
stretch of the run. A solver whose last accepted iterate lies above the target
is recorded as not reaching it at that $\varepsilon$. A run that dips below the
target, leaves it, and settles below it again is credited from the point at
which it settles.

Trajectories also expose \emph{plateaus}. The cost of a quadratically
convergent method takes very few distinct values across many decades of
$\varepsilon$: in \cref{tab:fwd}, Newton-SIN-BH costs 3, 4 or 5 \EVDs\ at
every target from $10^{-2}$ to $10^{-10}$. A single-tolerance comparison
therefore depends on where the tolerance falls relative to these plateaus.

\subsection{Accepted-only accounting}\label{sec:accounting}
A rejected line-search trial always costs one \EVD. Under
\emph{accepted-only} accounting it can never be credited with reaching a
target, because the method did not move to that point. Under
\emph{all-trial} accounting it can, which credits a method for an accuracy it
never returned. We use accepted-only accounting throughout and report the
other convention as a sensitivity check (\cref{sec:acct-results}); on the
family studied here the difference turns out to be small.

\subsection{Primitive work as the unit of cost}\label{sec:primitive}
\Cref{tab:primitive} gives the median primitive work at matched tolerance
($10^{-11}$) over the 270 instances.

\begin{table}[t]
\centering
\caption{Median primitive work at matched tolerance $10^{-11}$, 270 instances,
$n=100$. An `outer iteration' denotes a different amount of work in each
row.}
\label{tab:primitive}
\small
\begin{tabular}{lrrrr}
\toprule
solver & \EVDs & CG iterations & line-search trials & accepted outer its\\
\midrule
Newton-SIN-BH & 5     & 24   & 0  & 4\\
Newton-SIN    & 11    & 29.5 & 0  & 5\\
AGD-SDAJ-BH   & 70    & 0    & 14 & 9\\
AGD-SDAJ      & 81.5  & 0    & 25 & 10\\
SBB-Dual      & 46.5  & 0    & 0  & 45.5\\
Dykstra-APM   & 102.5 & 0    & 0  & 102.5\\
Anderson-APM  & 34    & 0    & 0  & 34\\
\bottomrule
\end{tabular}
\end{table}

Four accepted Newton iterations use 5 \EVDs\ (one at the initial point plus one
per accepted step) and 24 Krylov iterations. SBB-Dual,
Dykstra-APM and Anderson-APM use one \EVD\ per iteration and no further
primitive, although Anderson-APM's least-squares update carries $O(mn^2)$ work
that the \EVD\ count does not record. Nine AGD-SDAJ-BH outer
iterations use 70 \EVDs, because each outer iteration contains $q=2$
quasi-Newton inner steps and one accelerated gradient step, each with its own
line search, and the accelerated step can need a further evaluation
(\cref{alg:agd}). A table of
`iterations to convergence' across these rows would compare a different
quantity in each method class (Newton, AGD, SBB, APM). We therefore report \EVDs, Krylov iterations and
line-search trials separately.

\subsection{Timing protocol}\label{sec:timing}
All timings use Julia with \texttt{BLAS.set\_num\_threads(1)}. A warm-up solve
is run and discarded, \texttt{GC.gc()} is called before the measured solve,
and a single \texttt{@elapsed} surrounds the whole solve. A separate
instrumented pass (TimerOutputs) gives the breakdown of time. In that
breakdown, percentages are taken against the measured elapsed time, so any
uninstrumented work shows up as `other' time instead of inflating the
\EVD\ share. Timings were measured on all 270 instances at $n=100$, and on an
18-instance subset at $n=500$ (listed in \texttt{code/n500\_timing\_subset.txt}),
at both the native ($10^{-7}n$) and the matched ($10^{-11}$) tolerance. The
trajectory runs are a separate, untimed pass.

All seven variants were timed back to back in a single session, because
timings from different sessions are not comparable: an earlier session on the
same machine, with identical \EVD\ counts on every run, measured the dual
methods, whose code had not changed, between 3.7 and 4.0 times slower per
\EVD\ at $n=100$. Each
method is timed in an implementation that does not add work of its own. The
projection methods rebuild the projection from the positive eigenpairs, as the
dual methods do, rather than from all $n$, which removes an $O(n^3)$ product
per \EVD\ and gives bit-identical iterates on all 1104 projection-method runs
checked (270 $n=100$ instances and 6 $n=500$ instances, $\times$ 2 methods
$\times$ 2 tolerances; \texttt{code/validation/check\_projection\_rebuild.jl}); and
Anderson-APM is timed in a form that updates its buffers in place, with
iterates bit-identical to the transcription verified in \cref{app:anderson}.
Timed as transcribed, it would have cost $2.3$ to $3.8$ times more per \EVD.
All timings were measured on one machine: Intel Core 7 150U (10 physical
cores, 12 logical), 32\,GB RAM, Windows 11 (build 26200), Julia 1.12.6 with
OpenBLAS (ILP64) through libblastrampoline. The timing runs pin
\texttt{BLAS.set\_num\_threads(1)}; the untimed trajectory pass at $n=500$ was
run with four BLAS threads, which changes wall-clock only and not the \EVD\
counts it records.

\subsection{Fairness of the \EVD\ as a unit of work}\label{sec:evd-unit}
The \EVD\ count is a good proxy for cost only if one \EVD\ costs about the same
in every method. \Cref{tab:perevd} and \cref{fig:perevd} give the median
wall-clock time per \EVD\ at matched tolerance.

\begin{table}[t]
\centering
\caption{Microseconds per \EVD\ at matched tolerance ($10^{-11}$), $n=100$,
single-threaded BLAS.}
\label{tab:perevd}
\small
\begin{tabular}{lrr}
\toprule
solver & \si{\micro\second}/\EVD & relative to cheapest\\
\midrule
SBB-Dual      & 1133 & 1.00\\
AGD-SDAJ-BH   & 1181 & 1.04\\
AGD-SDAJ      & 1188 & 1.05\\
Dykstra-APM   & 1278 & 1.13\\
Anderson-APM  & 1323 & 1.17\\
Newton-SIN    & 1341 & 1.18\\
Newton-SIN-BH & 1695 & 1.50\\
\bottomrule
\end{tabular}
\end{table}

\begin{figure}[t]
\centering
\begin{tikzpicture}
\begin{axis}[
  width=0.8\linewidth, height=5.6cm,
  xbar, bar width=9pt,
  xmin=0, xmax=2000,
  xlabel={\si{\micro\second} per \EVD\ (matched tolerance $10^{-11}$)},
  ytick=data,
  % labels come from the data file, so they cannot fall out of step with the bars
  yticklabels from table={figures/data/per_evd_cost.dat}{name},
  y dir=reverse,
  nodes near coords, nodes near coords align={horizontal},
  every node near coord/.append style={font=\footnotesize},
  enlarge y limits=0.1,
  axis x line*=bottom, axis y line*=left,
]
\addplot[fill=gray!55, draw=gray!80!black] table[x=usperevd, y=idx] {figures/data/per_evd_cost.dat};
\end{axis}
\end{tikzpicture}
\caption{Median wall-clock time per \EVD\ at $n=100$, all seven variants timed
in one session. The spread is a factor of 1.50; Newton's \EVDs\ are the
dearest, because they are accompanied by Krylov solves that the \EVD\ count
does not record. Anderson-APM's \EVDs\ cost 1.17 times SBB-Dual's, the
difference coming from its work outside the eigendecomposition, chiefly its
least-squares update of $O(mn^2)$ work per \EVD.}
\label{fig:perevd}
\end{figure}


The median is taken over every run, including any that exhausted its budget
without converging (notably Newton-SIN's 56 naive-Armijo stalls of \cref{sec:trap2}).
In these timing tables a run's cost to its own exit counts as its cost ---
unlike the cost-to-target rule of \cref{sec:traj}, under which
a run that never reaches a target is recorded as not reaching it.
\Cref{tab:perevd} shows a factor-of-1.50 spread in per-\EVD\ cost across the
seven variants. Two solvers carry work the \EVD\ count does not see, and in
both cases counting \EVDs\ overstates that solver's advantage without
reversing it: Newton-SIN-BH is cheaper than SBB-Dual by 9.3 in \EVDs\
(5 against 46.5) but only 5.4 in wall-clock time (0.009\,s against
0.050\,s), because each Newton outer iteration hides an uncounted Krylov
solve; and Anderson-APM, whose least-squares update costs about 13--15\% of
its time against about 1\% for SBB-Dual, uses about 27\% fewer \EVDs\ than
SBB-Dual but only about 5\% less wall-clock time. In this single session, then,
the known bias in \EVD-counting is smaller than the gaps it is used to
compare. We report both metrics throughout, but for the reason given next we
treat \EVDs\ as primary and wall-clock time as a secondary check rather than
the reverse.

\paragraph{Session-to-session variation.}
\Cref{tab:perevd} is a single measured session. Repeating the
matched-tolerance timing measurement four further times shows a given
solver's median microseconds/\EVD\ varies far more across sessions (a
factor of 3.8 to 4.9, depending on the solver and dimension;
\texttt{results/analysis/timing\_repeats.md}) than the $\le\!1.50\times$
spread between solvers within one session --- a swing large enough, and
uniform enough across solvers within a session, that we read it as a
platform effect (\cref{sec:discussion}) rather than a property of the
methods. Two features nonetheless survive it: the cheapest solver per \EVD\
is always SBB-Dual or an AGD variant, at both dimensions, and the dearest per
\EVD\ is Newton-SIN-BH in 8 of the 10 sessions (Anderson-APM and Dykstra-APM
each take it once). It is the middle of the ranking that reorders, so the
specific ratios in \cref{tab:perevd} should be read as one session's draw,
not stable constants; the 5.4-in-time figure above is likewise approximate
to about ten per cent across sessions, while the
9.3-in-\EVDs\ figure, involving no timing measurement, is unaffected.
Ranking the seven variants by median \EVDs\ and by median wall-clock time
agrees in only 5 of these 10 sessions. The ranking by median \EVDs\ is
identical in all 10 sessions, since \EVD\ counts do not vary across sessions,
so every disagreement is wall-clock noise. By this second, wall-clock notion of
``dearest,'' Newton-SIN-BH is still cheapest overall in every session, but
the dearest \emph{overall} solver is Dykstra-APM in 8 of the 10 (a different
8 of 10 from the per-\EVD\ count above: all five at $n=500$, three of five
at $n=100$, where AGD-SDAJ overtakes it twice despite needing about a fifth
fewer \EVDs). So only the identity of the single cheapest method is
independent of the choice of cost metric; even within one session the
metric can decide a close pair (at $n=500$, Anderson-APM needs fewer \EVDs\
than AGD-SDAJ-BH on 13 of 18 instances but less time on only 11).

\subsection{Summary of the protocol}
\begin{enumerate}
  \item Measure forward error against an exact $\Xs$.
  \item Run every solver to a common tight tolerance and record a trajectory
    point after every \EVD.
  \item Read costs off the trajectories at a common target; give rejected
    trials no accuracy credit.
  \item Report \EVDs, Krylov iterations, line-search trials and wall-clock
    time, never iterations alone.
  \item Report the feasibility of the returned iterate
    ($\lambda_{\min}(X)$ and $\nrm{\diag(X)-e}_\infty$), and state which
    iterate is returned.
  \item Report reach rates beside every median that is computed over a
    subset of instances.
\end{enumerate}

\section{Two implementation traps}\label{sec:traps}

This section describes two implementation defects. Neither prevents
convergence to the correct solution at the tolerances usually reported, so
neither is visible from final errors. Both change measured cost substantially.
We present them as evidence for the protocol of \cref{sec:protocol}. They are
not new discoveries.

\subsection{Trap 1: redundant Jacobian factorization}\label{sec:trap1}
A semismooth Newton step at $y_k$ needs the spectral decomposition of
$C(y_k)$ to apply the generalized Jacobian~\cite{qisun2006quadratically}. That
decomposition was already computed when $\theta(y_k)$ and $\nabla\theta(y_k)$
were evaluated, at the accepted trial of the previous line search. An
implementation that recomputes it pays one extra \EVD\ per outer iteration and
gains nothing. On the literature matrix \texttt{cor1399} (\cref{sec:real}),
where both variants take the same six Newton steps, recomputing the
decomposition raises Newton's cost from 7 to 13 \EVDs, an increase of about
86\%, with an identical final error
(\texttt{results/timing\_session1/timing\_real\_matched.csv}). Newton-SIN-BH caches the decomposition throughout
this paper; Newton-SIN, the naive-Armijo control of \cref{sec:trap2}, does
not, so its \EVD\ counts already include this redundant factorization on top
of the naive test's own cost. (The 56-stall finding of \cref{sec:trap2} is
about convergence, not \EVD\ counts, so this confounding does not affect it.)
The same class of defect turned up in our first
reimplementation of AGD-SDAJ (\cref{app:agd}, defect~1). It was found in the
same way, by
counting primitives one at a time against a published count.

\subsection{Trap 2: the finite-precision Armijo test}\label{sec:trap2}
An objective-value sufficient-decrease test,
\begin{equation}\label{eq:armijo}
  \theta(y_k+\alpha d_k)\le\theta(y_k)+c\,\alpha\,\nabla\theta(y_k)^{\mathsf T}d_k,
\end{equation}
becomes a plain comparison of two rounded copies of $\theta_k$ once the
predicted decrease $c\,\alpha\,|\nabla\theta(y_k)^{\mathsf T}d_k|$ falls
below half an ulp of $\theta_k$: the right-hand side then rounds to
$\theta_k$ itself, and the test degenerates to $\mathrm{fl}(\theta_t)\le
\mathrm{fl}(\theta_k)$. It stops judging the step correctly once the
achieved decrease is itself buried in that same rounding error, near
$\theta_k$. Near the solution
$|\nabla\theta^{\mathsf T}d|\gtrsim\nrm{\nabla\theta}^2$ (equality for a
gradient-type step; for Newton's step it can be considerably larger), so
this second, more severe threshold --- which concerns whether the
\emph{achieved} decrease is resolvable at all, rather than the
$c\cdot\alpha$ sufficient-decrease margin --- falls at a gradient norm of roughly
$\sqrt{u\,|\theta|}$, where $u$ is the unit roundoff. Past that point,
whether a trial passes is decided by rounding noise rather than by the
step's real merit: a step that
improves both the gradient and the forward error can still be rejected, and
the search must halve until it happens to land where the noise lets a
trial through. In this regime Newton does not halt outright; it creeps,
accepting a string of steps far smaller than warranted until it exhausts its
iteration budget without converging --- a stall in the sense used below, as
the counts just below show. None of this is new. The
cancellation problem is well known in general
optimization~\cite{hager2005new}, and Higham~\cite[Ch.~17]{higham2002accuracy}
warns more broadly that stopping tests set below what rounding allows can be
unsatisfiable even by the rounded exact solution. For the NCM dual,
Borsdorf~\cite[\S4.7.1]{borsdorf2007newton} derived the threshold (about
$10^{-8}$ in $\nrm{\nabla\theta}$ when $|\theta|\approx1$) and proposed
alternatives, which became the
modified step selection of Borsdorf and Higham~\cite[Sec.~3.3,
Alg.~4.1]{borsdorf2010preconditioned}. Their Section~3.3 observes that when
$|\phi(x)|=1$, $\nrm{p}_2<u^{1/2}$ and $\nrm{\nabla\phi(x)}_2<u^{1/2}/2$, the
computed values $fl(\phi(x+p))$ and $fl(\phi(x))$ coincide, so that the Armijo
condition cannot be verified at all. Their Algorithm~4.1, Step~6, then tests
whether the two objective values are \emph{nearly equal}, using
$|a-b|<\gamma u(1+|a|+|b|)$ with $\gamma=100$, and if so accepts the full
inexact Newton step provided
$\nrm{\nabla f(y_k+d_k)}_2/\nrm{\nabla f(y_k)}_2\le1-\mu$, falling
back to a unit steepest-descent step otherwise. Our implementation uses their
parameter values, $\mubh=0.9$, $\rho=0.5$ and $\sigma=10^{-4}$ (we write
$\mubh$ for their gradient-reduction parameter, since $\mu$ denotes the
spectrum vector of \cref{def:family}).

\paragraph{Semismooth Newton.}
With the naive test~\eqref{eq:armijo}, Newton stalls (exits at
\texttt{max\_iter} without converging) on 56 of the 270 KKT
instances; across all 270 instances it uses 151{,}471 backtracks in total.
With Borsdorf--Higham globalization
(parameters $\mubh=0.9$, $\rho=0.5$, $\sigma=10^{-4}$, and $\gamma=100$ in the
near-equality test), it converges on all 270 with 8 backtracks in total. The count of
56 was reproduced exactly in a separate timing run of the same code: at
matched tolerance, the naive-Armijo control exits at \texttt{max\_iter} on
56/270. Both counts are in the shipped result files: the 56 is the
\texttt{exit} column of the \texttt{Newton-SIN} rows of
\texttt{timing\_kkt270\_matched.csv}, and the 8 is the
\texttt{linesearch\_trials} total of the \texttt{Newton-SIN-BH} rows of that
file, which agrees with the \texttt{backtracks} column of
\texttt{ranking\_kkt270\_v2.csv.gz}. The backtrack total of 151{,}471 is the
same run's own count of rejected trials, cross-checked against its
recorded per-trial trajectory (\texttt{code/trace\_armijo\_stall.jl}). Of the 5{,}677
rejected trials of the undamped Newton step ($\alpha=1$, before any
backtracking), every one, without exception, would, if accepted, have
reduced both the gradient norm and the exact forward error; counting every
rejected trial, including those reached only after backtracking, 115{,}211
of the 151{,}471 (76\%) would have.

\paragraph{AGD-SDAJ.}
The same failure occurs in AGD-SDAJ~\cite{huynh2025accelerated}, a separately
published method with a different algorithmic structure. Line~5 of
Algorithm~1 in that paper is the objective-value test~\eqref{eq:armijo}. Run
to the common tolerance $10^{-11}$, AGD-SDAJ as published exhausts the 4000-\EVD\
budget on 18 of the 270 instances (the budget is checked at the start of each
outer iteration, so a stalled run ends at between 4000 and 4035 \EVDs). On instance
\texttt{degen-r20-m1-d1e-10-p2}, the trajectory contains 3198 rejected
monotone line-search trials, with the forward error $\fro{X_k-\Xs}$ and
$\nrm{\nabla\theta}$ fixed at $6.0935\times10^{-9}$ and
$3.5986\times10^{-9}$. The predicted decrease is about $1.3\times10^{-21}$,
far below the rounding error in $\theta$, which is about $u|\theta|\approx
7\times10^{-14}$ because $\theta(\ys)\approx3\times10^{2}$ on this instance.
We adapted the Borsdorf--Higham near-equality test to
Algorithm~1's step: when the objective difference cannot be resolved, a trial
is accepted if it reduces $\nrm{\nabla\theta}$. The test is weaker than the
Newton safeguard's $(1-\mubh)$ reduction by design: that reduction is calibrated
to Newton steps, while near the precision floor a gradient step only inches the
gradient down, so the stronger test never fires. With
this change, all 18
stalls disappear (\cref{tab:agd-stall}). Across the family, the number of
instances reaching the tolerance rises from 252/270 to 270/270, and the
worst-case cost falls from 4035 \EVDs\ to 131.

\begin{table}[t]
\centering
\caption{AGD-SDAJ on \texttt{degen-r20-m1-d1e-10-p2}, as published and with
the Borsdorf--Higham adaptation.}
\label{tab:agd-stall}
\small
\begin{tabular}{lrrlrr}
\toprule
variant & outer & \EVDs & exit & final $\nrm{\nabla\theta}$ & $\fro{X-\Xs}$\\
\midrule
AGD-SDAJ (as published) & 115 & 4001 & \EVD\ budget & $3.60\times10^{-9}$ & $6.09\times10^{-9}$\\
AGD-SDAJ-BH             & 10  & 67   & converged    & $9.40\times10^{-12}$ & $1.82\times10^{-11}$\\
\bottomrule
\end{tabular}
\end{table}

\subsection{Scope of these findings}
Three points limit what these results show.
\begin{enumerate}
  \item \emph{The phenomenon is known.} Borsdorf analysed it for the NCM
    dual~\cite[\S4.7.1]{borsdorf2007newton}, and the fix was published
    in~\cite{borsdorf2010preconditioned}. What we add is limited to (i) its
    recurrence in a separately published method that uses a naive Armijo
    test, where it is invisible under that method's native stopping rule,
    and (ii) its quantification on the controlled family (56/270 and 18/270
    stalls). For semismooth Newton the remedy is the published
    Borsdorf--Higham modified line search; for AGD-SDAJ it is our adaptation
    of their near-equality safeguard to the Algorithm~1 step of that method.
  \item \emph{It does not affect Huynh and Hwang's reported results.} Their
    stopping rule, $\nrm{\nabla\theta}_2\le10^{-7}n$, equals $10^{-5}$ on
    these instances and stops the method far above the precision floor at
    which the test fails. The stall appears only when the method is pushed to
    a tight common tolerance, which is what a matched-accuracy comparison
    requires. On their own instances, our reimplementation reproduces their
    published counts exactly for P8 and P9, and matches P7 in outer
    iterations and residual with one differing backtracking decision
    (\cref{sec:repro}, \cref{app:agd}).
  \item \emph{The adaptation is ours.} AGD-SDAJ-BH is our modification and is
    not specified by the original authors. We therefore report it alongside
    the unmodified method on every test set where both were run (all except
    the 135 higher-rank instances of \cref{sec:highrank}).
\end{enumerate}
The general point is that whether a method appears to converge at a given
tolerance can depend on a line-search detail that is unrelated to the
algorithm's design. A benchmark that stops each method under its own rule
never exposes this.

\section{Ranking under different conventions}\label{sec:ranking}

\subsection{Solvers}\label{sec:solvers}
We compare five variants of four methods. All are implemented in Julia with
a common evaluation kernel, and every evaluation of $\theta$ or
$\nabla\theta$ costs exactly one \EVD. Unless stated otherwise, each variant
starts from $y_0=0$ (Dykstra-APM starts from $X_0=G$). The algorithms below
are written at the same level of detail and use the notation
$\bigl(\theta_k,g_k,X_k\bigr)\leftarrow\textsc{Eval}(y_k)$ for one \EVD\ of
$C(y_k)$, with $X_k=\Proj(C(y_k))$ and $g_k=\diag(X_k)-e$. The stopping
criterion in all algorithms is $\nrm{g_k}_2\le\tau$ (Dykstra-APM applies it to
the diagonal of its PSD iterate). Here $\tau=10^{-11}$ unless stated
otherwise.

\paragraph{Newton-SIN-BH} (\cref{alg:newton}) is the semismooth Newton
method of Qi and Sun~\cite{qisun2006quadratically}. The Newton system is solved inexactly by
conjugate gradients, using the block form of the generalized Jacobian
restricted to the positive-eigenvalue columns (with
$\Omega_{\beta\beta}=0$) and a $10^{-12}$ diagonal shift. The spectral
decomposition from the accepted trial is cached (\cref{sec:trap1}), and
globalization uses the Borsdorf--Higham modified line
search~\cite[Alg.~4.1, Step~6]{borsdorf2010preconditioned}. We adopt their
globalization only: their Algorithm~4.1 also solves the Newton system with
preconditioned \textsc{minres} rather than \textsc{cg}, which we do not
adopt, so \solver{Newton-SIN-BH} is not a reimplementation of their
algorithm. \emph{Newton-SIN} is the same method with the
naive Armijo test; it appears only as a control in
\cref{sec:protocol,sec:traps}.

\paragraph{AGD-SDAJ} (\cref{alg:agd}) is Algorithm~4 of Huynh and
Hwang~\cite{huynh2025accelerated}, transcribed from the full text and checked against
the published counts (\cref{app:agd}). It uses $q=2$ inner quasi-Newton steps
with a diagonal secant preconditioner and the nonmonotone residual line search
of their equation~(6), followed by one accelerated gradient step with a
monotone Armijo line search. \emph{AGD-SDAJ-BH} is our adaptation of that
Armijo step, described in \cref{sec:trap2}.

\paragraph{SBB-Dual} (\cref{alg:sbb}) is a spectral (Barzilai--Borwein)
gradient method on~\eqref{eq:dual}. It uses the BB1 step, clamped to
$[\eta,2-\eta]$ with $\eta=0.01$, has no line search, performs one \EVD\ per
update, and finishes with a diagonal-rescaling step
$\hat X=D^{-1/2}XD^{-1/2}$, $D=\Diag(\diag X)$, as in~\cite{borsdorf2010preconditioned}.

\paragraph{Dykstra-APM} (\cref{alg:apm}) is Higham's alternating projections
with Dykstra's correction~\cite{higham2002nearest}. Each iteration uses one \EVD. The
method returns the PSD half-iterate.

\begin{algorithm}[p]
\caption{Newton-SIN-BH (semismooth Newton, cached factorization, BH line search)}\label{alg:newton}
\begin{algorithmic}[1]
\Require $G$, tolerance $\tau$; $\sigma=10^{-4}$, $\mu=0.9$, $\gamma=100$, $u=2^{-53}$
\State $y\gets0$;\ $(\theta,g,X,\Lambda,P)\gets\textsc{Eval}(y)$ \Comment{keep the factorization}
\While{$\nrm{g}>\tau$}
  \State Solve $(V+10^{-12}I)d=-g$ by CG, with $V$ the generalized Jacobian from the cached $(\Lambda,P)$;
  \Statex \hspace{\algorithmicindent} stop at $\nrm{r}\le\min\{0.1,\nrm g\}\nrm g$ or on negative curvature
  \If{$d$ not finite or $g^{\mathsf T}d\ge0$} $d\gets-g$ \EndIf
  \State $\alpha\gets1$
  \Loop
    \State $(\theta_t,g_t,X_t,\Lambda_t,P_t)\gets\textsc{Eval}(y+\alpha d)$
    \If{$\theta_t\le\theta+\sigma\alpha g^{\mathsf T}d$} accept; \textbf{break} \EndIf
    \If{$|\theta_t-\theta|<\gamma u(1+|\theta_t|+|\theta|)$} \Comment{BH near-equality branch}
      \If{full step has $\nrm{g_{t}^{(\alpha=1)}}\le(1-\mu)\nrm g$} accept the full step; \textbf{break} \EndIf
      \State accept the unit steepest-descent step $y-g$ (one further \EVD); \textbf{break}
    \EndIf
    \State $\alpha\gets\alpha/2$
  \EndLoop
  \State update $(y,\theta,g,X,\Lambda,P)$ from the accepted trial \Comment{no recomputation}
\EndWhile
\State \Return $X$
\end{algorithmic}
\end{algorithm}

\begin{algorithm}[p]
\caption{AGD-SDAJ~\cite{huynh2025accelerated} (AGD-SDAJ-BH: with the marked branch)}\label{alg:agd}
\begin{algorithmic}[1]
\Require $G$, $\tau$; $q=2$, $c=10^{-4}$, $\rho=0.5$, $\lambda_1=\lambda_2=0.01$, $\epsilon_l=10^{-5}$, $\epsilon_u=10^{8}$
\State $y\gets0$;\ $(\theta,g,X)\gets\textsc{Eval}(y)$
\While{$\nrm{g}>\tau$}
  \State $(y^{p},\theta^{p},g^{p})\gets(y,\theta,g)$;\ $b\gets e$ \Comment{$B_0=I$, reset every outer iteration}
  \For{$i=0,\dots,q-1$ \textbf{while} $\nrm{g^p}>\tau$} \Comment{QN-SDAJ inner steps}
    \State $d\gets-g^p\oslash b$;\ $\epsilon_i\gets1/(i+1)^2$;\ $\alpha\gets1$
    \State backtrack $\alpha\gets\rho\alpha$ until, with $g_t=\nabla\theta(y^p+\alpha d)$,
    \Statex \hspace{2\algorithmicindent}$\nrm{g_t}^2\le(1+\epsilon_i)\nrm{g^p}^2-\lambda_1\alpha^2\nrm{g^p}^2-\lambda_2\alpha^2\nrm{d}^2$ \Comment{eq.~(6)}
    \State $\beta\gets a/b'$ with $a=-g^{p\mathsf T}d$, $b'=(g_t-g^p)^{\mathsf T}d$ if $b'>0$, else $\beta\gets1$ \Comment{general direction}
    \State $s\gets\beta\alpha d$;\ $y^+\gets y^p+s$;\ evaluate at $y^+$ (reuse the trial if $\beta=1$)
    \State $b_j\gets z_j/s_j$ if $z_j/s_j\in[\epsilon_l,\epsilon_u]$, else $1$, where $z=g^+-g^p$
    \State $(y^p,\theta^p,g^p)\gets(y^+,\theta^+,g^+)$
  \EndFor
  \If{$\theta^p>\theta$} $(y^p,\theta^p,g^p)\gets(y,\theta,g)$ \EndIf
  \If{$\nrm{g^p}\le\tau$} $y\gets y^p$; \textbf{break} \Comment{early exit at $x^{\mathrm{pre}}$}
  \EndIf
  \State $d\gets-g^p$;\ $\alpha\gets1$
  \Loop \Comment{Algorithm 1, line 5}
    \State $(\theta_t,g_t)\gets\textsc{Eval}(y^p+\alpha d)$
    \If{$\theta_t\le\theta^p+c\alpha g^{p\mathsf T}d$} \textbf{break} \EndIf
    \If{\textsc{bh} \textbf{and} $|\theta_t-\theta^p|<\gamma u(1+|\theta_t|+|\theta^p|)$ \textbf{and} $\nrm{g_t}<\nrm{g^p}$} \textbf{break} \EndIf
    \State $\alpha\gets\rho\alpha$
  \EndLoop
  \State $\beta\gets\nrm{g^p}^2/\bigl((g^p-g_t)^{\mathsf T}g^p\bigr)$ if the denominator is $>0$, else $1$
  \State $y\gets y^p-\beta\alpha g^p$;\ $(\theta,g,X)\gets\textsc{Eval}(y)$ (reuse the trial if $\beta=1$)
\EndWhile
\State \Return $X$
\end{algorithmic}
\end{algorithm}

\begin{algorithm}[t]
\caption{SBB-Dual (clamped spectral gradient on the dual)}\label{alg:sbb}
\begin{algorithmic}[1]
\Require $G$, $\tau$, $\eta=0.01$
\State $y\gets0$;\ $(\theta,g,X)\gets\textsc{Eval}(y)$
\While{$\nrm{g}>\tau$}
  \If{first update, or $s^{\mathsf T}w\le0$, where $s=y-y_{\mathrm{prev}}$, $w=g-g_{\mathrm{prev}}$}
    \State $t\gets1$
  \Else
    \State $t\gets\min\{\max\{s^{\mathsf T}s/s^{\mathsf T}w,\ \eta\},\ 2-\eta\}$ \Comment{clamped BB1}
  \EndIf
  \State $(y_{\mathrm{prev}},g_{\mathrm{prev}})\gets(y,g)$;\ $y\gets y-tg$;\ $(\theta,g,X)\gets\textsc{Eval}(y)$
\EndWhile
\State $X\gets D^{-1/2}XD^{-1/2}$, $D=\Diag(\diag X)$;\ set $\diag(X)=e$ \Comment{rescaling epilogue}
\State \Return $X$
\end{algorithmic}
\end{algorithm}

\begin{algorithm}[t]
\caption{Dykstra-APM (alternating projections with Dykstra correction)}\label{alg:apm}
\begin{algorithmic}[1]
\Require $G$, $\tau$
\State $X\gets G$;\ $\Delta S\gets0$;\ $\Delta U\gets0$
\Loop
  \State $Y\gets\Proj(X+\Delta S)$ \Comment{one \EVD}
  \If{$\nrm{\diag(Y)-e}_2\le\tau$} \Return $Y$ \Comment{PSD half-iterate}
  \EndIf
  \State $\Delta S\gets X+\Delta S-Y$
  \State $Z\gets Y+\Delta U$ with $\diag(Z)\gets e$;\ $\Delta U\gets Y+\Delta U-Z$;\ $X\gets Z$
\EndLoop
\end{algorithmic}
\end{algorithm}

\subsection{Design}
All five variants were run on the 270 instances of \cref{sec:family} at
$n=100$, to a common tolerance of $10^{-11}$ with a budget of 4000 \EVDs, and
a trajectory was recorded as described in \cref{sec:traj}. We evaluated 11
conventions: the native exit; a common forward-error target
$\fro{X-\Xs}\le\varepsilon$ for
$\varepsilon\in\{10^{-2},10^{-4},10^{-6},10^{-8},10^{-10}\}$; and a common
dual residual $\nrm{\nabla\theta}_2\le\tau$ for the same five values. We also
compared accepted-only with all-trial accounting.
\Cref{sec:n500} repeats the study at $n=500$.

\subsection{The winner does not change}
\Cref{tab:conventions} lists the orderings produced by the 11 conventions.
There are two distinct orderings, and Newton-SIN-BH is cheapest in both.
\emph{This is a null result for winner reversal.} On this family, the
conventions change cost ratios and the order of the lower places. They never
change which method is cheapest.

\begin{table}[t]
\centering
\caption{Distinct orderings of median cost across the 11 conventions.}
\label{tab:conventions}
\small
\begin{tabular}{p{0.52\linewidth}p{0.4\linewidth}}
\toprule
ordering (cheapest first) & conventions\\
\midrule
Newton-SIN-BH $<$ SBB-Dual $<$ AGD-SDAJ-BH $<$ AGD-SDAJ $<$ Dykstra-APM &
forward $10^{-6},10^{-8},10^{-10}$; residual $10^{-6},10^{-8},10^{-10}$; native\\
Newton-SIN-BH $<$ AGD-SDAJ-BH $<$ AGD-SDAJ $<$ SBB-Dual $<$ Dykstra-APM &
forward $10^{-2},10^{-4}$; residual $10^{-2},10^{-4}$\\
\bottomrule
\end{tabular}
\end{table}

\subsection{The ordering below the winner changes with the accuracy target}
\Cref{tab:fwd,tab:res} and \cref{fig:crossover} show median cost to a common
target. At loose targets ($10^{-2}$, $10^{-4}$), both AGD variants are cheaper
than SBB-Dual. From $10^{-6}$ onward, SBB-Dual is cheaper. The same reversal
occurs at the same point under the forward-error and the dual-residual
criteria, so it does not depend on which error measure is used. A paper that
reports at a loose tolerance would place AGD-SDAJ above SBB-Dual, and a paper
that reports at a tight tolerance would place them the other way. Both are
defensible one-number summaries of the same runs.

\begin{figure}[t]
\centering
\begin{tikzpicture}
\begin{axis}[
  width=0.9\linewidth, height=7cm,
  xmode=log, x dir=reverse,
  xmin=5e-11, xmax=2e-1,
  ymin=0, ymax=100,
  xlabel={target forward error $\varepsilon$ in $\fro{X-\Xs}\le\varepsilon$},
  ylabel={median \EVDs\ to reach and hold target},
  legend pos=north west, legend cell align=left,
  legend style={font=\footnotesize},
  grid=major, grid style={gray!20},
]
\addplot[black, thick, mark=*, mark size=1.3pt] table[x=eps, y=NewtonSINBH] {figures/data/evd_vs_eps.dat};
\addlegendentry{Newton-SIN-BH}
\addplot[blue!70!black, thick, mark=square*, mark size=1.3pt] table[x=eps, y=SBBDual] {figures/data/evd_vs_eps.dat};
\addlegendentry{SBB-Dual}
\addplot[red!70!black, thick, mark=triangle*, mark size=1.6pt] table[x=eps, y=AGDSDAJ] {figures/data/evd_vs_eps.dat};
\addlegendentry{AGD-SDAJ$^\dagger$}
\addplot[orange!85!black, thick, dashed, mark=triangle, mark size=1.6pt] table[x=eps, y=AGDSDAJBH] {figures/data/evd_vs_eps.dat};
\addlegendentry{AGD-SDAJ-BH$^\dagger$}
\addplot[green!45!black, thick, mark=diamond*, mark size=1.6pt] table[x=eps, y=DykstraAPM] {figures/data/evd_vs_eps.dat};
\addlegendentry{Dykstra-APM}
\addplot[violet, thick, mark=pentagon*, mark size=1.6pt] table[x=eps, y=AndersonAPM] {figures/data/evd_vs_eps.dat};
\addlegendentry{Anderson-APM}
\end{axis}
\end{tikzpicture}
\caption{Median cumulative \EVDs\ to reach and hold $\fro{X_k-\Xs}\le\varepsilon$
(accepted-only accounting) over 270 instances, $n=100$, on a half-decade grid.
Anderson-APM is second at every target shown. SBB-Dual and AGD-SDAJ-BH cross
between $\varepsilon\approx3\times10^{-5}$ and $10^{-5}$. $^\dagger$AGD-SDAJ-BH and every other solver reach and hold each
target shown on all 270 instances. Unmodified AGD-SDAJ does so on all 270 down
to $3\times10^{-7}$, and on 253/270 at $10^{-10}$; below $3\times10^{-7}$ its
medians are taken over the instances it reached. The curves stop at
$10^{-10}$.
Cluster-bootstrap intervals for the decade targets are given in
\cref{tab:fwd}; they are narrower than the plotted marks at most targets.}
\label{fig:crossover}
\end{figure}


\begin{table}[t]
\centering
\caption{Median \EVDs\ to reach and hold $\fro{X-\Xs}\le\varepsilon$
(accepted-only). Entries marked $^{*}$ are medians over the subset of
instances that reached the target; reach counts are in brackets. All
unmarked entries are over 270/270.}
\label{tab:fwd}
\begin{tabular}{lrrrrr}
\toprule
$\varepsilon$ & Newton-SIN-BH & AGD-SDAJ-BH & AGD-SDAJ & SBB-Dual & Dykstra-APM\\
\midrule
$10^{-2}$  & \textbf{3} & 6 & 6 & 8 & 15\\
$10^{-4}$  & \textbf{4} & 14 & 14 & 16 & 33.5\\
$10^{-6}$  & \textbf{4} & 28 & 28 & 24.5 & 53.5\\
$10^{-8}$  & \textbf{5} & 45 & 50$^{*}$ [264] & 34 & 74.5\\
$10^{-10}$ & \textbf{5} & 63.5 & 72$^{*}$ [253] & 43.5 & 95.5\\
\bottomrule
\end{tabular}
\end{table}

\begin{table}[t]
\centering
\caption{Median \EVDs\ to reach and hold $\nrm{\nabla\theta}_2\le\tau$
(accepted-only), with reach counts as in \cref{tab:fwd}.}
\label{tab:res}
\begin{tabular}{lrrrrr}
\toprule
$\tau$ & Newton-SIN-BH & AGD-SDAJ-BH & AGD-SDAJ & SBB-Dual & Dykstra-APM\\
\midrule
$10^{-2}$  & \textbf{2} & 5 & 5 & 7 & 12\\
$10^{-4}$  & \textbf{3} & 13 & 13 & 14 & 30\\
$10^{-6}$  & \textbf{4} & 27 & 27 & 23 & 50\\
$10^{-8}$  & \textbf{5} & 42 & 47$^{*}$ [264] & 32.5 & 71\\
$10^{-10}$ & \textbf{5} & 60 & 69$^{*}$ [254] & 41.5 & 92\\
\bottomrule
\end{tabular}
\end{table}

\paragraph{Reach and hold.}
A target counts as reached at the first accepted iterate from which the error
stays at or below the target for the rest of the run. The accelerated methods
are not monotone in forward error, so this differs from the first iterate at
which the error dips below the target; scoring the first dip alone would
understate their cost. Under this rule Newton-SIN-BH, AGD-SDAJ-BH, SBB-Dual
and Dykstra-APM reach every target on all 270 instances. Only unmodified
AGD-SDAJ misses some: 264/270 and 253/270 at forward targets $10^{-8}$ and
$10^{-10}$ (264 and 254 for the residual), which is consistent with its 18
budget-exhausting stalls (\cref{sec:trap2}). Its entries marked $^{*}$ are
medians over the instances it reached, and are therefore optimistic. The
second/third-place reversal is the same whether or not those entries are
included.
\paragraph{Uncertainty.}
\Cref{tab:fwdci} gives cluster-bootstrap 95\% intervals for the forward-error
medians ($B=2000$, resampling paired seeds within rank, so that all 18
$(m,\delta)$ cells of a drawn seed move together). The intervals are narrow
and do not overlap between SBB-Dual and the AGD variants at any target, so the
change of order between $10^{-4}$ and $10^{-6}$ is not a sampling artifact.

\begin{table}[t]
\centering
\caption{Cluster-bootstrap 95\% intervals for the medians of \cref{tab:fwd}
($n=100$, $B=2000$, resampling the five paired seeds within each rank).
Bracketed counts are instances reaching and holding the target.}
\label{tab:fwdci}
\small
\begin{tabular}{lrrrrr}
\toprule
$\varepsilon$ & Newton-SIN-BH & AGD-SDAJ-BH & AGD-SDAJ & SBB-Dual & Dykstra-APM\\
\midrule
$10^{-2}$  & 3 [3, 3] & 6 [6, 6.5] & 6 [6, 6.5] & 8 [8, 8] & 15 [15, 15]\\
$10^{-4}$  & 4 [3, 4] & 14 [13, 15] & 14 [13, 15] & 16 [15.5, 16] & 33.5 [33.5, 33.5]\\
$10^{-6}$  & 4 [4, 4] & 28 [27, 29] & 28 [27, 29] & 24.5 [24.5, 24.5] & 53.5 [53.5, 53.5]\\
$10^{-8}$  & 5 [5, 5] & 45 [42, 48.5] & 50 [46, 52] \tiny[264] & 34 [34, 34] & 74.5 [74.5, 74.5]\\
$10^{-10}$ & 5 [5, 5] & 63.5 [62, 66] & 72 [70, 75] \tiny[253] & 43.5 [43, 43.5] & 95.5 [95.5, 95.5]\\
\bottomrule
\end{tabular}
\end{table}

The paired per-instance picture at $n=100$ is, however, much closer than the
aggregate medians suggest. Against AGD-SDAJ-BH at $\varepsilon=10^{-2}$ the
median paired difference is only $+2$ \EVDs\ (95\% CI $[1,2]$), SBB-Dual is
cheaper on 40.4\% of the 270 instances, and 10 instances tie. The $n=100$
ordering between these two is therefore a modest average effect with
substantial per-instance overlap, in contrast to $n=500$ (\cref{sec:n500}),
where SBB-Dual is cheaper on 0 of 270 instances at the same target.

\subsection{The ordering varies across the degeneracy grid}\label{sec:cells}
Aggregate medians can hide dependence on the regime. At a fixed target
(forward error $10^{-8}$), the 54 $(r,m,\delta)$ cells produce 7 distinct
orderings (\cref{tab:cells}).

\begin{table}[t]
\centering
\caption{Orderings by cell at forward target $10^{-8}$ (54 cells, five paired
instances each), with the number of cells of each rank producing each
ordering.}
\label{tab:cells}
\small
\begin{tabular}{p{0.56\linewidth}rrrr}
\toprule
ordering (cheapest first) & cells & $r{=}5$ & $r{=}20$ & $r{=}50$\\
\midrule
Newton-SIN-BH $<$ SBB-Dual $<$ AGD-SDAJ-BH $<$ AGD-SDAJ $<$ Dykstra-APM & 19 & 0 & 13 & 6\\
Newton-SIN-BH $<$ AGD-SDAJ-BH $<$ AGD-SDAJ $<$ SBB-Dual $<$ Dykstra-APM & 18 & 17 & 1 & 0\\
Newton-SIN-BH $<$ SBB-Dual $<$ Dykstra-APM $<$ AGD-SDAJ-BH $<$ AGD-SDAJ & 13 & 0 & 2 & 11\\
Newton-SIN-BH $<$ AGD-SDAJ-BH $<$ SBB-Dual $<$ AGD-SDAJ $<$ Dykstra-APM & 1 & 0 & 1 & 0\\
Newton-SIN-BH $<$ SBB-Dual $<$ AGD-SDAJ-BH $<$ Dykstra-APM $<$ AGD-SDAJ & 1 & 0 & 1 & 0\\
Newton-SIN-BH $<$ AGD-SDAJ $<$ AGD-SDAJ-BH $<$ SBB-Dual $<$ Dykstra-APM & 1 & 1 & 0 & 0\\
Newton-SIN-BH $<$ SBB-Dual $<$ Dykstra-APM $<$ AGD-SDAJ $<$ AGD-SDAJ-BH & 1 & 0 & 0 & 1\\
\bottomrule
\end{tabular}
\end{table}

Newton-SIN-BH is cheapest in every cell. Below it, at $n=100$, the order is
governed almost entirely by the rank $r$ of the solution. AGD-SDAJ-BH is
cheaper than SBB-Dual in all 18 cells with $r=5$; SBB-Dual is cheaper in 16 of
the 18 cells with $r=20$ and in all 18 with $r=50$. Within a rank, the
near-zero multiplicity $m$ and the separation $\delta$ rarely change
this pair's order. Both AGD variants fall below Dykstra-APM to last place in
14 cells: all 12 cells with $r=50$ and $m\in\{5,20\}$, and the two cells with
$r=20$, $m=20$ and $\delta\in\{10^{-8},10^{-10}\}$. This block is determined
by rank and multiplicity; the separation plays no visible role in it. The
instance family was built to vary these quantities one at a time. On a
collection of real matrices with uncontrolled spectra, the ordering could not
be associated with rank in this way.

The attribution is nevertheless specific to $n=100$. \Cref{sec:n500} repeats
the study at $n=500$, where SBB-Dual is no longer cheaper at any rank. What
survives across both dimensions is the direction of the effect, not the
crossover: SBB-Dual gains on the AGD variants as rank increases.

That rankings depend on the required accuracy is expected in
principle~\cite{more2009benchmarking,hansen2021coco}. What the controlled
family adds is that the variation is organized by a single spectral quantity
which the construction sets directly.

Each cell contains only five instances, so a per-cell median is a noisy
statistic, and some of the rarer orderings may reflect sampling variation
rather than a real change of regime.
Resampling the five paired seeds of each cell ($B=2000$) shows that the count
of distinct orderings overstates how much genuinely varies. Most cells
reproduce their observed ordering in nearly every resample, but the adjacent
pair AGD-SDAJ-BH against AGD-SDAJ is usually \emph{unresolved}: the 95\%
interval of its median paired difference contains zero in most cells, which is
unsurprising because the two differ only in a line-search safeguard that rarely
activates at $n=100$. Reading unresolved adjacent pairs as ties collapses much
of the variety in \cref{tab:cells}. What survives is the placement of SBB-Dual
relative to the AGD pair, which is resolved in the great majority of cells and
is the comparison discussed above.
\Cref{fig:heatmap} shows the same comparison cell by cell as a ratio of median
costs, which separates the three controlled quantities: rank sets the level,
the near-zero multiplicity modulates it within a rank, and the separation has
almost no effect.

\begin{figure}[t]
\centering
\begin{tikzpicture}
\begin{axis}[
  width=0.86\linewidth, height=7.4cm,
  colormap/viridis,
  colorbar,
  colorbar style={
    title={$\;$median \EVDs: SBB-Dual $/$ AGD-SDAJ-BH},
    title style={font=\footnotesize, rotate=90, xshift=2.1cm, yshift=-2.4cm},
    ylabel style={font=\footnotesize},
  },
  point meta min=0.3, point meta max=2.8,
  xlabel={separation $\rho=\delta/\mu_{\mathrm{bulk}}$},
  ylabel={rank $r$ and near-zero multiplicity $m$},
  xtick={0,1,2,3,4,5},
  xticklabels={$0$,$10^{-9}$,$10^{-7}$,$10^{-5}$,$10^{-3}$,$10^{-1}$},
  ytick={0,1,2,3,4,5,6,7,8},
  yticklabels={%
    {$r{=}5,m{=}1$},{$r{=}5,m{=}5$},{$r{=}5,m{=}20$},%
    {$r{=}20,m{=}1$},{$r{=}20,m{=}5$},{$r{=}20,m{=}20$},%
    {$r{=}50,m{=}1$},{$r{=}50,m{=}5$},{$r{=}50,m{=}20$}},
  yticklabel style={font=\footnotesize},
  xticklabel style={font=\footnotesize},
  tick style={draw=none},
  enlargelimits=false,
  axis on top,
]
\addplot[matrix plot*, point meta=explicit, mesh/cols=6]
  table[x=x, y=y, meta=ratio] {figures/data/heatmap.dat};
\end{axis}
\end{tikzpicture}
\caption{Ratio of median \EVDs\ to reach and hold $\fro{X-\Xs}\le10^{-8}$,
SBB-Dual divided by AGD-SDAJ-BH, for each of the 54 cells at $n=100$ (five
paired instances per cell). Values above 1 mean AGD-SDAJ-BH is cheaper; below
1, SBB-Dual is. The ratio is organized primarily by rank: it exceeds 1 in
every $r=5$ cell and is below 1 in every $r=50$ cell. The near-zero
multiplicity $m$ modulates the level within a rank, most visibly at $r=20$
(from $1.03$ at $m=1$ to $0.40$ at $m=20$), whereas the separation $\rho$ has
little effect, as the near-constancy along each row shows.}
\label{fig:heatmap}
\end{figure}


\subsection{The same study at $n=500$}\label{sec:n500}
The instance family was regenerated at $n=500$ with the same grid. All 270
instances screened successfully; the acceptance rate at $r=5$ fell to 5 of 7
candidate seeds, the first rank to show rejections, which is what
\cref{prop:sat} predicts. The ranking study uses all five paired seeds of
each rank, 270 instances, with caps of 1200 iterations for SBB-Dual and 800
for Dykstra-APM. It was run in two parts, seeds 0--2 (about 8.6\,h) and seeds
3--4 (about 7.1\,h), with identical code and settings.

\begin{table}[t]
\centering
\caption{$n=500$, 270 instances. Median \EVDs\ to reach and hold
$\fro{X-\Xs}\le\varepsilon$, with cluster-bootstrap 95\% intervals
($B=2000$, resampling paired seeds within rank) and reach counts where these
are below 270.}
\label{tab:n500}
\small
\begin{tabular}{lrrrrr}
\toprule
$\varepsilon$ & Newton-SIN-BH & AGD-SDAJ-BH & AGD-SDAJ & SBB-Dual & Dykstra-APM$^{\ddagger}$\\
\midrule
$10^{-2}$  & \textbf{3} [3, 3] & 9 [9, 9] & 9 [9, 9] & 44 [44, 44] & 87 [86.5, 88]\\
$10^{-4}$  & \textbf{4} [3, 4] & 30 [29, 32] & 30 [29, 32] & 97 [97, 97] & 195 [194.5, 195.5]\\
$10^{-6}$  & \textbf{4} [4, 4] & 53.5 [50.5, 56] & 55.5 [52, 58] \tiny[268] & 152 [152, 152] & 191 [191, 192] \tiny[180]\\
$10^{-8}$  & \textbf{5} [5, 5] & 79 [76, 80] & 87 [84, 89] \tiny[259] & 207.5 [207, 208] & 258 [258, 260] \tiny[180]\\
$10^{-10}$ & \textbf{5} [5, 5] & 101.5 [98.5, 105] & 119 [117, 123] \tiny[242] & 260 [259, 262] \tiny[261] & 325.5 [325, 327.5] \tiny[180]\\
\bottomrule
\end{tabular}
\end{table}

Two of the $n=100$ conclusions survive and two do not.

\emph{Confirmed.} Newton-SIN-BH is again cheapest at every target, at 5 \EVDs.
The stall of unmodified AGD-SDAJ recurs (31 of 270 runs exhaust the \EVD\
budget, against none for AGD-SDAJ-BH, whose worst case is 240 \EVDs\ against
4039).

\emph{Not reproduced.} The change of order with the accuracy target seen in
\cref{tab:fwd} is absent: at $n=500$ a single ordering holds at all five
targets, with AGD-SDAJ-BH ahead of SBB-Dual throughout. The paired
per-instance comparison is one-sided. SBB-Dual is cheaper on 0 of 270
instances at $10^{-2}$, 3 at $10^{-4}$ and 15 at $10^{-8}$. We report
these counts and the cluster-bootstrap intervals of \cref{tab:n500} rather
than an instance-level significance test: the design reuses each drawn $U$
across all 18 $(m,\delta)$ cells of its rank, so the 270 comparisons form 15
seed clusters, and a test treating them as independent would overstate
the evidence by a wide margin. The rank effect of \cref{sec:cells} keeps its
direction but no longer produces a crossover: the median paired difference at
$10^{-8}$ falls from $+669.5$ \EVDs\ at $r=5$ to $+129$ at $r=20$ and $+13$ at
$r=50$, where SBB-Dual is cheaper on 17\% of instances. Cell by cell, SBB-Dual
is ahead of AGD-SDAJ-BH in 0 of the 54 cells with a resolved bootstrap
interval (47 resolved in favour of AGD-SDAJ-BH, 7 unresolved). SBB-Dual is thus
behind at every rank tested at $n=500$, while it led at $r\ge20$ at $n=100$.

This is not an artifact of the caps. SBB-Dual converged on all 270 instances
(worst case 1075 \EVDs\ against its 1200 cap) and reached every target except
$10^{-10}$ on nine instances. $^{\ddagger}$Dykstra-APM, by contrast, \emph{is}
censored: 90 of 270 runs stop at the 800-iteration cap, and from $10^{-6}$
down its medians cover only 180 instances. Its position at $n=500$ should be
read as censored rather than measured.

The scaling is the cleaner way to state the difference. Between $n=100$ and
$n=500$ the median cost to $10^{-8}$ rises by a factor of 6.1 for SBB-Dual
(34 to 207.5 \EVDs) and by 1.8 for AGD-SDAJ-BH (45 to 79), while Newton-SIN-BH
is unchanged at 5. The order below the winner therefore reflects the test set,
including its dimension, rather than a fixed property of the methods --- which
is the paper's point, in a stronger form than the $n=100$ data alone supports.

\subsection{Four matrices from the literature}\label{sec:real}
The family of \cref{sec:family} is synthetic. To check that the regime
dependence above is not an artifact of the construction, we ran the same
protocol on four invalid correlation matrices from the literature: the
Rocky Mountain carbon-storage matrix ($n=94$), \texttt{cor1399} ($n=1399$),
\texttt{cor3120} ($n=3120$) and \texttt{bccd16} ($n=3250$), all from the
Anymatrix \texttt{CORRINV} collection~\cite{higham2022anymatrix} used
by~\cite{huynh2025accelerated}.

No exact solution is available here, which is the situation the rest of this
paper avoids. Forward error is measured against a reference computed by
semismooth Newton driven towards $10^{-13}$; for \texttt{cor1399} that
reference solve itself stopped at its iteration limit after 1770 \EVDs\ with
$\nrm{\nabla\theta}_2=4.59\times10^{-13}$. At $n\approx3000$ an absolute
tolerance of $10^{-13}$ is not attainable in double precision, so for
\texttt{cor3120} and \texttt{bccd16} the reference solve is capped at 30
\EVDs\ and its achieved residual is reported instead:
$\nrm{\nabla\theta}_2=1.58\times10^{-12}$ and $8.07\times10^{-13}$, after
2201\,s and 895\,s. We therefore treat these runs as a check on external
validity, not as measurements of the same quality as \cref{tab:fwd}. The common
target is $10^{-9}$; each solver is capped at 500 \EVDs\ on the two smaller
matrices and at 300 iterations on the two larger ones. Wall-clock time was
measured only on the two smaller matrices.

\begin{table}[t]
\centering
\caption{Real matrices, target $10^{-9}$. \EVDs\ to reach and hold $10^{-8}$,
total \EVDs, measured time at the matched tolerance (smaller two matrices
only) and final error. $^{\ddagger}$Stops at its cap and is censored.}
\label{tab:real}
\small
\begin{tabular}{llrrrr}
\toprule
matrix & solver & \EVDs\ to $10^{-8}$ & total \EVDs & time (s) & final error\\
\midrule
Rocky Mountain & Newton-SIN-BH & \textbf{4} & 4 & 0.009 & $1.49\times10^{-13}$\\
$n=94$         & SBB-Dual      & \textbf{7} & 8 & 0.011 & $1.53\times10^{-10}$\\
               & Dykstra-APM   & 11 & 13 & 0.023 & $2.57\times10^{-10}$\\
               & AGD-SDAJ-BH   & 12 & 17 & 0.023 & $3.47\times10^{-10}$\\
               & AGD-SDAJ      & 12 & 25 & 0.036 & $3.47\times10^{-10}$\\
\midrule
\texttt{cor1399} & Newton-SIN-BH & \textbf{7} & 7 & 7.00 & $1.27\times10^{-9}$\\
$n=1399$         & AGD-SDAJ-BH   & \textbf{66} & 91 & 65.00 & $6.92\times10^{-10}$\\
                 & AGD-SDAJ      & 67 & 136 & 100.04 & $2.29\times10^{-9}$\\
                 & SBB-Dual      & 200 & 212 & 152.69 & $1.89\times10^{-9}$\\
                 & Dykstra-APM$^{\ddagger}$ & 406 & 300 & 288.90 & $1.49\times10^{-6}$\\
\midrule
\texttt{cor3120} & Newton-SIN-BH & \textbf{6} & 6 & -- & $2.33\times10^{-9}$\\
$n=3120$         & AGD-SDAJ-BH   & 41 & 42 & -- & $1.32\times10^{-9}$\\
                 & AGD-SDAJ      & 56 & 59 & -- & $2.17\times10^{-10}$\\
                 & SBB-Dual      & 255 & 270 & -- & $3.38\times10^{-9}$\\
                 & Dykstra-APM$^{\ddagger}$ & -- & 300 & -- & $1.16\times10^{-5}$\\
\midrule
\texttt{bccd16}  & Newton-SIN-BH & \textbf{4} & 5 & -- & $2.71\times10^{-11}$\\
$n=3250$         & SBB-Dual      & 5 & 6 & -- & $1.20\times10^{-10}$\\
                 & Dykstra-APM   & 6 & 7 & -- & $1.46\times10^{-10}$\\
                 & AGD-SDAJ-BH   & 8 & 9 & -- & $8.88\times10^{-10}$\\
                 & AGD-SDAJ$^{\ddagger}$ & -- & 323 & -- & $3.93\times10^{-5}$\\
\bottomrule
\end{tabular}
\end{table}

The matrices disagree with each other, broadly in the direction the
synthetic family predicts. At $n=94$, SBB-Dual is the cheapest of the
first-order methods, 7 \EVDs\ to $10^{-8}$ against 12 for AGD-SDAJ-BH and 11
for Dykstra-APM, which is the pattern of \cref{tab:cells} at $n=100$. At
$n=1399$, SBB-Dual costs three times AGD-SDAJ-BH (200 against 66), and at
$n=3120$ six times (255 against 41), which is the pattern of \cref{sec:n500}.
\texttt{bccd16} is easy for every method: all but one reach $10^{-8}$ within
8 \EVDs, so it does not discriminate between the first-order methods.
Newton-SIN-BH is cheapest on all four, as everywhere else in this paper. The
change of order is therefore visible on matrices with uncontrolled spectra that
we did not construct, and on the two smaller matrices the same ordering holds
under wall-clock time. With four matrices that differ in more than dimension,
this is consistent with the dimension effect of \cref{sec:n500}, not an
independent demonstration of it.

\paragraph{The line-search stall occurs on real data.}
On the Rocky Mountain matrix the naive-Armijo Newton control exhausts its
budget at a forward error of $1.03\times10^{-8}$, while the same solver with
the Borsdorf--Higham line search converges in 4 \EVDs\ to
$1.49\times10^{-13}$. On \texttt{cor1399} the naive control converges normally
(13 \EVDs\ against 7). On \texttt{bccd16} the same defect appears in
unmodified AGD-SDAJ. Both AGD variants follow an identical path to a forward
error of $3.93\times10^{-5}$ at the fifth \EVD. From there AGD-SDAJ-BH, whose
safeguard accepts the step, converges in four more \EVDs\ to
$8.9\times10^{-10}$; unmodified AGD-SDAJ rejects trial after trial in its
monotone line search, with the error unchanged in its first eight digits, until
it exhausts the budget at 323 \EVDs. The defect of \cref{sec:trap2} is thus not
an artifact of the constructed family, and it affects both published methods on
real data, but neither is it universal: whether it appears depends on the
matrix.

\paragraph{The cost of a reference.}
At $n\approx3000$ the reference solve alone took 15 to 37 minutes, and a first
attempt at \texttt{bccd16} with an uncapped absolute tolerance ran for four
hours without finishing. A reference accurate enough to measure forward error
can cost more than the solves being measured. That is a practical argument for
the synthetic family, whose solution is available in closed form at any size.

\subsection{Accounting convention}\label{sec:acct-results}
Under all-trial accounting, every evaluated point counts toward reaching and
holding a target, including rejected line-search trials. A rejected trial can
then be credited with reaching a target, but one that lies above the target
also breaks a hold. On this family the second effect dominates slightly.
Median costs change only for the two AGD variants (\cref{tab:acct}), by at
most 2 \EVDs, and all-trial accounting is marginally the more expensive.
Newton-SIN-BH, SBB-Dual and Dykstra-APM are unaffected, because they have few
or no rejected trials here. The effect is negligible on this family and never
changes an ordering. For methods with many rejected trials it need not be,
and accepted-only accounting is the rule that cannot credit a point the method
itself rejected.

\begin{table}[t]
\centering
\caption{Accepted-only vs all-trial accounting (forward error, reach and
hold). Only rows where the median changed are shown; bracketed counts are
instances reaching the target.}
\label{tab:acct}
\begin{tabular}{llrrr}
\toprule
$\varepsilon$ & solver & accepted-only & all-trial & $\Delta$\\
\midrule
$10^{-4}$  & AGD-SDAJ-BH & 14.0 & 15.0 & $+1.0$\\
$10^{-4}$  & AGD-SDAJ    & 14.0 & 15.0 & $+1.0$\\
$10^{-8}$  & AGD-SDAJ-BH & 45.0 & 46.5 & $+1.5$\\
$10^{-8}$  & AGD-SDAJ    & 50.0 [264] & 52.0 [264] & $+2.0$\\
$10^{-10}$ & AGD-SDAJ-BH & 63.5 & 65.5 & $+2.0$\\
$10^{-10}$ & AGD-SDAJ    & 72.0 [253] & 74.0 [253] & $+2.0$\\
\bottomrule
\end{tabular}
\end{table}

\subsection{What each solver returns}\label{sec:returned}
Methods also differ in what they return (\cref{tab:feas}). All five return
matrices with $\lambda_{\min}$ at roundoff level, so none is meaningfully more
positive semidefinite than another. The diagonal is a different matter.
SBB-Dual returns an exactly unit diagonal because of its rescaling epilogue,
whereas every other variant returns a diagonal that is accurate only to its
convergence tolerance. The worst-case diagonal error of unmodified AGD-SDAJ
is $4.01\times10^{-8}$, about 2.5 orders of magnitude larger than that of any
other variant, which is a direct consequence of the stalls in
\cref{sec:trap2}. The rescaling epilogue has a cost of its own: on a
nearly feasible input at loose tolerance, it can increase $\fro{X-\Xs}$.
Borsdorf and Higham bound that increase by
$\mathrm{tol}\,\fro{X}/(1-\mathrm{tol})$~\cite[Lem.~3.2]{borsdorf2010preconditioned},
which is small whenever the tolerance is tight but need not be at a loose
one.
A benchmark that reports only cost hides these differences. A user who needs
an exactly unit diagonal and a user who needs minimal distance may not be
best served by the same method.

\begin{table}[t]
\centering
\caption{Feasibility of the returned iterate at matched tolerance $10^{-11}$,
270 instances. ``diag err'' is $\nrm{\diag(X)-e}_\infty$.}
\label{tab:feas}
\begin{tabular}{lrrrr}
\toprule
solver & median $\lambda_{\min}(X)$ & worst $\lambda_{\min}$ & median diag err & worst diag err\\
\midrule
Newton-SIN-BH & $-3.42\times10^{-15}$ & $-1.25\times10^{-14}$ & $8.38\times10^{-15}$ & $5.70\times10^{-12}$\\
AGD-SDAJ-BH   & $-3.40\times10^{-15}$ & $-1.36\times10^{-14}$ & $2.53\times10^{-12}$ & $7.20\times10^{-12}$\\
AGD-SDAJ      & $-3.40\times10^{-15}$ & $-1.39\times10^{-14}$ & $2.79\times10^{-12}$ & $\mathbf{4.01\times10^{-8}}$\\
SBB-Dual      & $-3.42\times10^{-15}$ & $-1.49\times10^{-14}$ & $0$ & $0$\\
Dykstra-APM   & $-3.43\times10^{-15}$ & $-1.36\times10^{-14}$ & $1.49\times10^{-12}$ & $9.37\times10^{-12}$\\
\bottomrule
\end{tabular}
\end{table}

\subsection{Native exits and exit reasons}
For completeness, \cref{tab:nativeexit} reports cost at each solver's own
exit when all are run to $10^{-11}$, and \cref{tab:exits} the exit reasons.
Under this convention, the ordering equals the tight-target ordering in
\cref{tab:conventions}, because the common tolerance is tight.

\begin{table}[t]
\centering
\caption{Cost at the exit of each solver under the common tolerance $10^{-11}$.}
\label{tab:nativeexit}
\begin{tabular}{lrrrr}
\toprule
solver & median \EVDs & min & max & median final $\fro{X-\Xs}$\\
\midrule
Newton-SIN-BH & 5   & 4  & 10   & $1.25\times10^{-13}$\\
AGD-SDAJ-BH   & 70  & 27 & 131  & $1.10\times10^{-11}$\\
AGD-SDAJ      & 82  & 34 & 4035 & $1.16\times10^{-11}$\\
SBB-Dual      & 46  & 13 & 218  & $1.74\times10^{-11}$\\
Dykstra-APM   & 102 & 24 & 443  & $1.98\times10^{-11}$\\
\bottomrule
\end{tabular}
\end{table}

\begin{table}[t]
\centering
\caption{Exit reasons, 270 instances. ``at $x^{\mathrm{pre}}$'' denotes the
early exit at the preconditioned point of~\cite{huynh2025accelerated}.}
\label{tab:exits}
\begin{tabular}{llr}
\toprule
solver & exit & count\\
\midrule
Newton-SIN-BH & converged & 270\\
AGD-SDAJ-BH   & tolerance met / at $x^{\mathrm{pre}}$ & 177 / 93\\
AGD-SDAJ      & tolerance met / at $x^{\mathrm{pre}}$ / \EVD\ budget & 156 / 96 / 18\\
SBB-Dual      & tolerance met & 270\\
Dykstra-APM   & tolerance met & 270\\
\bottomrule
\end{tabular}
\end{table}

\section{The dimension-scaled stopping rule}\label{sec:stopping}

A common exit criterion for dual NCM methods is
$\nrm{\nabla\theta(y_k)}_2\le10^{-7}n$. It is used, for example, in the
experiments of~\cite{huynh2025accelerated}. Because the threshold grows linearly
with $n$, the forward error it guarantees depends on both the dimension and
the method. The general point that native stopping tests confound
comparisons was made by Dolan, Mor\'e and Munson~\cite{dolan2006optimality}.
This section documents a specific NCM instance of it.

\paragraph{Spread across methods at a fixed $n$.}
\Cref{tab:native} and \cref{fig:native} report the timing run under this rule
at $n=100$, where the tolerance is $10^{-5}$.

\begin{table}[t]
\centering
\caption{Solvers exiting when the stopping quantity specific to each method
first reaches the common numerical threshold $10^{-7}n$ at $n=100$ (for the
dual methods this quantity is $\nrm{\nabla\theta}_2$; for Dykstra-APM it is
the diagonal error of the PSD iterate), 270 instances, single-threaded
timing.}
\label{tab:native}
\begin{tabular}{lrrr}
\toprule
solver & median \EVDs & median time (s) & median $\fro{X-\Xs}$\\
\midrule
Newton-SIN-BH & 4  & 0.023 & $5.04\times10^{-8}$\\
Newton-SIN    & 7  & 0.035 & $5.04\times10^{-8}$\\
AGD-SDAJ-BH   & 20 & 0.085 & $9.58\times10^{-6}$\\
AGD-SDAJ      & 20 & 0.084 & $9.58\times10^{-6}$\\
SBB-Dual      & 18 & 0.088 & $7.73\times10^{-6}$\\
Dykstra-APM   & 40 & 0.222 & $1.93\times10^{-5}$\\
\bottomrule
\end{tabular}
\end{table}

\begin{figure}[t]
\centering
\begin{tikzpicture}
\begin{axis}[
  width=0.8\linewidth, height=5.4cm,
  xbar, bar width=9pt,
  xmode=log, xmin=1e-8, xmax=1e-4, log origin=infty,
  xlabel={median $\fro{X-\Xs}$ at the method-specific exit at threshold $10^{-7}n$},
  ytick={1,...,6},
  yticklabels={Newton-SIN-BH,Newton-SIN,AGD-SDAJ-BH,AGD-SDAJ,SBB-Dual,Dykstra-APM},
  y dir=reverse,
  enlarge y limits=0.12,
  axis x line*=bottom, axis y line*=left,
  grid=major, grid style={gray!20},
]
\addplot[fill=gray!55, draw=gray!80!black] table[x=err, y=idx] {figures/data/native_accuracy.dat};
\end{axis}
\end{tikzpicture}
\caption{Achieved accuracy when each solver exits under the same published
dimension-scaled rule (tolerance $10^{-5}$ at $n=100$). Median forward
errors range from $5.04\times10^{-8}$ (Newton variants) to
$1.93\times10^{-5}$ (Dykstra-APM), a factor of 383. Costs measured at these
exits are therefore costs at different accuracies. Anderson-APM was added after
these timing runs and is not shown.}
\label{fig:native}
\end{figure}


At this dimension the best and worst median achieved accuracies differ by a
factor of 383. That factor is specific to $n=100$: repeating the same run at
$n=500$ (18 instances, threshold $5\times10^{-5}$) gives a spread of only
13.8, with medians from $1.74\times10^{-5}$ for Newton-SIN-BH to
$2.41\times10^{-4}$ for Dykstra-APM. The spread across solvers is therefore
not itself a stable quantity, and we do not claim it as one. What does grow
with $n$ is the inaccuracy the rule permits, which is the subject of the next
paragraph. Newton
overshoots the tolerance because a quadratically convergent step moves far
past the threshold. First-order methods stop close to it, and Dykstra-APM's
stopping quantity (the diagonal error of its PSD iterate) is not the dual
gradient norm at all. The cost column is therefore a comparison at different
accuracies, and it favours methods that stop early. At this tolerance, it also
ranks SBB-Dual (18 \EVDs) ahead of both AGD variants (20 \EVDs), which is the
opposite of the matched-accuracy ordering at the corresponding forward error
(\cref{tab:fwd}, $\varepsilon=10^{-4}$ to $10^{-6}$).

\paragraph{Growth with $n$.}
The rule loosens as $n$ grows. In earlier runs, outside the ranking study, a
nearly feasible $n=500$ instance exited after \emph{zero} updates with forward
error $7.8\times10^{-4}$. At $n=2000$, both dual first-order methods stopped
after 3 to 5 iterations with forward errors between $2.7\times10^{-3}$ and
$4.1\times10^{-3}$. Under this rule, a larger problem can appear easier
simply because less accuracy is being asked for.
The paired family makes the same point with the dimension controlled. Under
this rule Newton-SIN-BH exits at a median forward error of
$5.04\times10^{-8}$ at $n=100$ but $1.74\times10^{-5}$ at $n=500$, about 340
times worse, purely because the threshold is proportional to $n$. Every other
variant degrades similarly. A reader comparing published $n=100$ and $n=500$
results under this rule is comparing solutions of very different quality.
The real-matrix runs of \cref{sec:real} were made at a matched tolerance
only, so they give no native-rule spread on real data.

\paragraph{Recommendation.}
Timing or iteration counts should never be reported without the achieved
forward error (or, when $\Xs$ is unknown, a certified bound on it). Where a
dimension-scaled rule is used for compatibility with earlier work, it should
be accompanied by a run to a common tight tolerance.

\section{Reproducibility}\label{sec:repro}

\paragraph{Independent implementations of the earlier solvers.}
SBB-Dual, semismooth Newton and Dykstra-APM were implemented independently in
Python and in Julia, using \emph{different} constructions of the reference
solution. On all 14 shared instances, the two implementations agree on update
counts and forward errors to every printed digit.

\paragraph{AGD-SDAJ: verified against the published counts.}
AGD-SDAJ exists only in Julia. Because this paper argues that implementation
details can reverse conclusions, we did not report any result for our
reimplementation on the KKT family until it had reproduced Huynh and Hwang's
own published counts on their own instances. Those instances are three
matrices from their test set, loaded under the same convention used elsewhere
in this project (\cref{tab:agd-verify}).

\begin{table}[t]
\centering
\caption{AGD-SDAJ ($q=2$): published counts (outer iterations / \EVDs\ /
nonmonotone backtracks / $\fro{A-\Xs}$) against our reimplementation.
Published values are from Table~3 and Table~A1 of~\cite{huynh2025accelerated}.}
\label{tab:agd-verify}
\begin{tabular}{llll}
\toprule
instance & published & this work & status\\
\midrule
P9 \texttt{bccd16}, $n=3250$  & 1 / 5 / 0 / 29.06  & 1 / 5 / 0 / 29.0563 & exact\\
P8 \texttt{cor3120}, $n=3120$ & 3 / 18 / 1 / 5.44  & 3 / 18 / 1 / 5.4437 & exact\\
P7 \texttt{cor1399}, $n=1399$ & 4 / 31 / 8 / 21.03 & 4 / 32 / 7 / 21.0339 & one backtrack differs\\
\bottomrule
\end{tabular}
\end{table}

P9 and P8 agree exactly. P7 agrees on outer iterations and residual norm, and
differs by one backtracking decision: 7 rejected trials locally against 8
published. That single decision accounts for the whole difference in \EVDs.
Excluding backtracks, the local base cost is $32-7=25$ and the published base
cost is $31-8=23$. The gap of 2 \EVDs\ is exactly the early-termination saving
described in the note to their Table~3: with the extra backtrack, their
trajectory reaches the tolerance \emph{at} the preconditioned point and skips
the final gradient step, whereas ours reaches it one step later. We attribute
the differing decision to floating-point differences. The reference
implementation is in MATLAB and ours in Julia. The two call different LAPACK
eigensolver drivers, and on a dense $1399\times1399$ problem, differences in
the last bits of the spectrum are enough to flip a single acceptance test
whose two sides are nearly equal. We did not search over parameters to make
P7 match, because tuning to a target would be exactly the kind of hidden
fitting this paper warns against. \Cref{app:agd} gives the full verification
record, including three transcription defects found and corrected along the
way.

\paragraph{Single-machine caveat.}
All timings come from one machine with single-threaded BLAS: Intel Core 7
150U (10 physical cores, 12 logical), 32\,GB RAM, Windows 11 (build 26200),
Julia 1.12.6 with OpenBLAS (ILP64) through libblastrampoline. Per-\EVD\ ratios
(\cref{tab:perevd}) may differ on other hardware and BLAS libraries, and with
multithreading.
Each configuration was timed once, after a discarded warm-up, so we do not
report run-to-run variance.

\paragraph{Artifact.}
The generator, the instances with hashes, all solver code, the trajectory
and timing CSV files, and the analysis scripts that produce every table in
this paper are listed in \cref{app:artifact}.
They are available at \url{https://github.com/Realosunboy6/ncm-exact-benchmark}. An archived
snapshot carrying a permanent identifier will be deposited on acceptance.

\section{Discussion and limitations}\label{sec:discussion}

\paragraph{What changed the conclusions, and what did not.}
On this family, none of the conventions we tested changed which method was
cheapest. The accuracy target, the instance regime and the problem dimension
did change the order below it: the change of lower-place ordering with the
accuracy target that we found at $n=100$ does not reproduce at $n=500$
(\cref{sec:n500}), and the rank effect that organizes it at $n=100$ no longer
changes the order there. The stopping rule
changed the achieved accuracy by a factor of 383. The cost unit (\EVDs\ against wall-clock time) and the accounting rule
(accepted-only against all-trial) changed cost ratios but no ordering. Two
line-search details changed whether methods converged at all. We think the
correct summary is conservative. The protocol matters mainly for the
\emph{magnitude} of reported advantages and for the ordering of methods that
are close in cost. On the evidence here, it does not matter for identifying a
clear winner. We expect this to hold wherever one method is ahead by a
large factor; where methods are within a factor of about two, conclusions
drawn under a single convention should be treated with caution.

\paragraph{Why an exact family.}
The regime dependence in \cref{sec:cells} is organized mainly by rank, one of
the three quantities the construction controls. Finding it requires both an
exact $\Xs$, so that forward error is not contaminated by an approximate
reference, and independent control of the spectral structure at the
solution. Real matrices supply neither. The same machinery then showed the
limit of that finding: at $n=500$ (\cref{sec:n500}) the rank effect keeps its
direction but no longer changes the order, so what looked like a rule about
the methods at one dimension turns out to depend on the benchmark regime as
well.

\paragraph{Limitations.}
\begin{itemize}
  \item \emph{Problem size.} Results are reported at $n=100$ and $n=500$
    (270 instances each). The $r\ge20$ crossover between SBB-Dual and
    AGD-SDAJ found at $n=100$ does \emph{not} persist at $n=500$, where
    AGD-SDAJ-BH is cheaper at every rank. No dimension above 500 was tested
    on the family. Dykstra-APM is censored at $n=500$ by its iteration cap.
  \item \emph{Synthetic family.} The KKT family is synthetic. Its $G$
    matrices are correlation-like but are not estimated from data.
    \Cref{sec:real} adds four matrices from the literature, from $n=94$ to
    $n=3250$, measured against a computed rather than an exact reference; for
    the two largest that reference is capped and reaches
    $\nrm{\nabla\theta}_2\approx10^{-12}$ rather than $10^{-13}$, and no
    wall-clock times were taken. Four matrices cannot separate dimension from
    the other ways in which they differ.
  \item \emph{Statistical uncertainty.} Aggregate medians at both dimensions
    carry cluster-bootstrap intervals (\cref{tab:fwdci,tab:n500}). Each
    dimension rests on five paired seeds per rank, that is 15 seed clusters,
    so inferential statements remain modest. Cell-level
    orderings rest on five instances each, and within a cell the adjacent pair
    AGD-SDAJ-BH against AGD-SDAJ is usually unresolved, so the count of
    distinct orderings should not be read as a count of real regime changes.
    We report counts and intervals rather than significance tests, because the
    reuse of each drawn $U$ across all cells of its rank makes instance-level
    comparisons strongly dependent.
  \item \emph{Single machine.} All timings come from one machine with
    single-threaded BLAS.
  \item \emph{Our adaptation.} The Borsdorf--Higham adaptation to AGD-SDAJ is
    ours, not the original authors'. It is reported beside the unmodified
    method throughout.
  \item \emph{Feasibility propositions.} \Cref{prop:det} assumes uniform
    $\mu$, and the extension to mixed $\mu$ is open (\cref{app:family}).
    \Cref{prop:hd} is stated for $a$ only; no companion bound on $b$ is
    given, since the released family lies in the fixed-rank regime where
    \cref{prop:hd} does not apply.
  \item \emph{Scope of methods.} We benchmark five variants of four methods.
    Other methods (e.g.\ Anderson-accelerated alternating
    projections~\cite{higham2016anderson}) are not included.
    Neither are the recent semismooth Newton method for general projection
    equations~\cite{armijo2025semismooth} or the factor-structured
    variant~\cite{borsdorf2010factor}.
  \item \emph{Attribution of the line search.} We use the Borsdorf--Higham
    globalization but not the rest of their algorithm, and the adaptation of
    their near-equality safeguard to AGD-SDAJ is ours
    (\cref{sec:solvers,sec:trap2}).
  \item \emph{Unweighted problem.} Only the unweighted Frobenius problem is
    considered.
\end{itemize}

\paragraph{Conclusion.}
An exact instance family with a prescribed solution structure makes it
possible to separate stopping-rule, reference, accounting and
primitive-work effects in NCM benchmarks. On 270 instances at $n=100$, these
effects changed cost ratios and the lower places of the ranking, but not the
cheapest method. Repeating the study on 270 instances at $n=500$ sharpened
that finding: Newton-SIN-BH remains the leading method, but the change of
lower-place ordering with the accuracy target does not reproduce, and the rank
effect that organizes it at $n=100$ no longer changes the order. Rankings
below the leading method are therefore not intrinsic properties of the
solvers; they depend materially on the benchmark regime, including its
dimension. The one clear failure we found is a known finite-precision
line-search defect, which appears in two separately published methods, on
the constructed family and on real matrices. It is
removed in the semismooth Newton method by the published Borsdorf--Higham
modified line search, and in AGD-SDAJ by our adaptation of their near-equality
safeguard. We recommend
that NCM comparisons report forward error against a trusted solution at a
common tolerance, give rejected trials no accuracy credit, and report
primitive work and wall-clock time together.


\appendix
\section{Proofs for the instance family}\label{app:family}

Throughout, $U\in\mathbb{R}^{n\times r}$ has rows $u_i$, $\Xs=UU^{\mathsf T}$,
$P_r$ is the orthogonal projector onto $\range U=\range\Xs$, and
$Q_0Q_0^{\mathsf T}=I-P_r$ (because $\Xs$ is symmetric, $\ker\Xs$ is the
orthogonal complement of $\range\Xs$).

\subsection{Proof of \cref{prop:det}}
With $\mu=c\,e$, we have $W=cQ_0Q_0^{\mathsf T}=c(I-P_r)$. For $i\ne j$, the
diagonal correction $\Diag(\diag W)$ in~\eqref{eq:construction} does not
affect the entry $(i,j)$, and $I_{ij}=0$, so
\[
  G_{ij}=\Xs_{ij}-W_{ij}=\Xs_{ij}+c\,(P_r)_{ij}.
\]
By the triangle inequality, $|G_{ij}|\le|\Xs_{ij}|+c|(P_r)_{ij}|\le a+cb$. If
$a<1$ and $c\le(1-a)/b$, then $a+cb\le1$. \qed

\begin{remark}
The bound is sufficient but not necessary. Cancellation between $\Xs_{ij}$ and
$c(P_r)_{ij}$, and the fact that $a$ and $b$ are generally attained at
different pairs $(i,j)$, both make it conservative. For this reason the
generator checks $|G_{ij}|\le1$ entry by entry.
\end{remark}

\subsection{Proof of \cref{prop:sat}}
Fix $r\ge2$ and $\epsilon\in(0,1)$. Because $S^{r-1}$ is compact, it can be
covered by finitely many caps
$C_1,\dots,C_{N}$, $N=N(\epsilon)$, each of angular radius $\epsilon/2$. Any
two points in the same cap make an angle of at most $\epsilon$, so their inner
product is at least $\cos\epsilon\ge1-\epsilon^2/2$. Once $n>N$, the
pigeonhole principle forces two of $u_1,\dots,u_n$ into a common cap, so
$a\ge1-\epsilon^2/2$ deterministically. Since $\epsilon$ is arbitrary and
$a\le1$ always, $a\to1$ along every realization. (A probabilistic argument is
needed only for a \emph{rate}; the limit itself follows from pigeonhole.)
\qed

\begin{remark}[What this does not give]
It is tempting to conclude that a fixed $c>0$ must eventually violate
\cref{prop:det}. That does not follow. At fixed $r$ the projector bound $b$
also tends to zero, so $a+cb\le1$ can persist: with $a_n=1-2/n$ and
$b_n=1/n$ one has $a_n\to1$ while $a_n+cb_n<1$ for every $n$ at $c=1$. A claim
in either direction needs the relative rates of $1-a$ and $b$, which we do not
establish here.
\end{remark}

\subsection{Proof of \cref{prop:hd}}
Fix $i\ne j$ and condition on $u_i$. Since $u_j$ is uniform on $S^{r-1}$ and
independent of $u_i$, the uniform distribution on the sphere is subgaussian:
by~\cite[Thm.~3.4.5]{vershynin2026high}, for every unit vector $v$ and every
$t\ge0$,
\[
  \Pr\bigl\{u_j^{\mathsf T}v\ge t\bigr\}\le 2\exp(-t^2r/2).
\]
Applying this with $v=u_i$, and again to $-u_j$, gives
$\Pr\{|u_i^{\mathsf T}u_j|\ge t\}\le4\exp(-t^2r/2)$; the bound is unchanged
after averaging over $u_i$. Taking a union bound over the fewer than $n^2/2$
unordered pairs,
\[
  \Pr\{a\ge t\}\;\le\;2n^2\exp(-t^2r/2).
\]
Setting the right-hand side equal to $\delta$ and solving for $t$ gives
$t=\sqrt{(2/r)(2\log n+\log(2/\delta))}$, which is the stated bound. If
$r/\log n\to\infty$ then $t\to0$ for any fixed $\delta$. \qed

\begin{remark}[On $b$]
We do not state a companion bound for $b=\max_{i\ne j}|(P_r)_{ij}|$. Writing
$P_r=VV^{\mathsf T}$ with $V$ having orthonormal columns, $(P_r)_{ij}$ is an
inner product of two rows of $V$, whose joint distribution is not that of two
independent uniform vectors; a bound would need a model for $\range U$ (for
instance a Haar-random subspace) together with a justification that the
row-normalized Gaussian $U$ induces it. Since the released family lies in the
fixed-rank regime of \cref{prop:sat}, where \cref{prop:hd} does not apply, and
since the generator screens $|G_{ij}|\le1$ entrywise in any case, we leave
this open.
\end{remark}

\begin{remark}[Scope]
These propositions concern $|G_{ij}|\le1$ only. Indefiniteness of $G$
($\lambda_{\min}(G)<0$) is not implied and is checked separately. The
released family ($r\in\{5,20,50\}$, $n=100$) is in the fixed-rank regime of
\cref{prop:sat}, not the high-dimensional regime of \cref{prop:hd}.
\end{remark}

\subsection{Mixed $\mu$}
For the released $\mu=(\mu_{\mathrm{bulk}}e_{k-m},\ \delta e_m)$, write
$Q_0=[Q_b\ Q_\delta]$. Then
$W=\mu_{\mathrm{bulk}}Q_bQ_b^{\mathsf T}+\delta Q_\delta Q_\delta^{\mathsf T}
=\mu_{\mathrm{bulk}}(I-P_r)-(\mu_{\mathrm{bulk}}-\delta)Q_\delta Q_\delta^{\mathsf T}$.
The off-diagonal entries therefore pick up an extra term
$(\mu_{\mathrm{bulk}}-\delta)(Q_\delta Q_\delta^{\mathsf T})_{ij}$, whose size
depends on the choice of basis for $\ker\Xs$ used to select $Q_\delta$. A
bound analogous to \cref{prop:det} follows with $b$ replaced by
$b+\max_{i\ne j}|(Q_\delta Q_\delta^{\mathsf T})_{ij}|$.

\begin{remark}[Open]
We do not characterize $\max_{i\ne j}|(Q_\delta Q_\delta^{\mathsf T})_{ij}|$.
It depends on which orthonormal basis of $\ker\Xs$ is chosen, and the
generator fixes that basis through the eigendecomposition it computes, so the
quantity is not determined by $(n,r,m,\delta)$ alone. A basis-free bound would
need either a canonical choice of $Q_\delta$ or a bound uniform over
admissible bases. We therefore state \cref{prop:det} for uniform $\mu$ only,
and the generator screens the mixed-$\mu$ instances entrywise rather than
relying on a bound. This is why the released family is screened rather than
certified.
\end{remark}

\section{AGD-SDAJ verification record}\label{app:agd}

\paragraph{Gate.}
Before running our reimplementation of AGD-SDAJ~\cite{huynh2025accelerated} on the
KKT family, we required it to reproduce the authors' own published counts on
their own instances. An unverified reimplementation that is then reported as
losing would not be a defensible result.

\paragraph{Source and implementation.}
The source is Huynh and Hwang (2025), \emph{Numer.\ Linear Algebra Appl.}\
32(6), DOI \texttt{10.1002/nla.70045}. We transcribed Algorithm~1 (AGD),
Algorithm~3 (QN-SDAJ), Algorithm~4 (AGD-SDAJ) and the Section~4.3 parameters
from the full text: $c=10^{-4}$, $\rho=0.5$, $\lambda_1=\lambda_2=0.01$,
$\epsilon_l=10^{-5}$, $\epsilon_u=10^{8}$, $q=2$, $x_0=0$, the stopping rule
$\nrm{\nabla\theta}\le10^{-7}n$, and $k_{\max}=200$. The implementation is
\texttt{agd\_sdaj.jl} (standalone, used for verification) and
\texttt{agd\_sdaj\_ncm} in \texttt{bench\_sbb\_dual.jl} (instrumented, used for
the study). The instances are the Anymatrix \texttt{CORRINV} files already used
elsewhere in this project, exported with the same loading convention.

\paragraph{Equation (6), as transcribed.}
The nonmonotone line search of Algorithm~3 is a Li--Fukushima-type
rule~\cite{li2000derivative} on the residual:
\begin{equation}\label{eq:hh6}
  \nrm{\nabla\theta(x+\alpha d)}^2\ \le\ (1+\epsilon_k)\nrm{\nabla\theta(x)}^2
  -\lambda_1\nrm{\alpha\nabla\theta(x)}^2-\lambda_2\nrm{\alpha d}^2,
  \qquad \epsilon_k=\frac{1}{(k+1)^2},
\end{equation}
where, as implemented, $k$ is the index of the inner step within its outer
iteration ($\epsilon_i$ in Algorithm~\ref{alg:agd}), so that
$\epsilon_k\in\{1,\tfrac14\}$ for $q=2$. This indexing is our reading of the
text, not stated explicitly there; with it, the published counts of
\cref{tab:agd-full} are reproduced exactly on P8 and P9. The test uses
$\lambda_1=\lambda_2=0.01$, and $\alpha=\rho^{\,j}$ is backtracked from 1 over
$j=0,1,2,\dots$ (we write $j$ for the backtracking exponent, since $m$ is the
near-zero multiplicity of Definition~\ref{def:family}). The
merit function is $\tfrac12\nrm{\nabla\theta}^2$. The quasi-Newton direction
$d=-B^{-1}\nabla\theta$ need not be a descent direction for it, which is why a
nonmonotone rule is needed. The paper also confirms that $B_k$ is reset to the
identity once $x_k^{\mathrm{pre}}$ is generated, and that Step~4 discards
$x_k^{\mathrm{pre}}$ if it increases $\theta$.

Before the full text was available, we tried two guesses: a
Grippo--Lampariello--Lucidi rule on $\theta$, and a Li--Fukushima rule. The
second had the right general form but three wrong details: unsquared norms, a
missing $(1+\epsilon_k)$ slack, and $\lambda_1,\lambda_2$ attached to the
wrong terms. The guess columns of \cref{tab:agd-full} were run before the
acceleration-factor correction described below, so they are not directly
comparable with the final column. On P8 their \EVD\ counts bracketed the
published one; on P7 both were well above it, which pointed to a remaining
defect outside the line search (item~3 below).

\paragraph{Results.}
\begin{table}[h]
\centering
\caption{Outer iterations / \EVDs\ / nonmonotone backtracks, and
$\fro{G-\Xs}$, under the final transcription (\texttt{eq6}) and the two
superseded guesses (run at an earlier code state; see text).}
\label{tab:agd-full}
\small
\begin{tabular}{lllll}
\toprule
instance & published & \texttt{eq6} (final) & GLL guess & LF guess\\
\midrule
P9 \texttt{bccd16}  & 1/5/0, 29.06  & 1/5/0, 29.0563 (exact) & 1/5/0 & 1/5/0\\
P8 \texttt{cor3120} & 3/18/1, 5.44  & 3/18/1, 5.4437 (exact) & 3/17/0 & 3/19/2\\
P7 \texttt{cor1399} & 4/31/8, 21.03 & 4/32/7, 21.0339 & 6/40/3 & 5/48/17\\
\bottomrule
\end{tabular}
\end{table}
Under the final transcription, $\fro{G-\Xs}$ matched the published value to
the displayed precision on all three instances. This showed that the
objective, gradient, projection and stopping rule were correct, and that any
discrepancy lay in the amount of work per iteration.

\paragraph{Three hazards when reimplementing Algorithm~4.}
Each of the following is easy to introduce when transcribing the method from
its pseudocode. Each leaves the method converging to the correct solution
while changing the work it does, so none is visible from final errors, only
from work counts checked against the published ones. We record them for anyone
reimplementing the method. All three occurred in our own first transcription
and were found by exactly that check: the first two by accounting for the five
published \EVDs\ on P9 one at a time, the third because a gap in outer
iterations on P7 remained once the line search was correct.
\begin{enumerate}
  \item \emph{Redundant \EVD\ at an already-evaluated point.} The inner
    QN-SDAJ routine recomputed $\theta$ and $\nabla\theta$ at the point its
    caller had just evaluated. This is the same class of defect as
    \cref{sec:trap1}, and it was found in the same way. Fix: pass the
    caller's values into the inner routine.
  \item \emph{Missing early termination at the preconditioned point.} The
    note to Table~3 of~\cite{huynh2025accelerated} states that the loop also stops
    whenever $x_k^{\mathrm{pre}}$ satisfies the stopping criterion, which
    saves at least one \EVD\ (the trial at $\alpha=1$), and two when
    $\beta\ne1$. Without this check, the outer gradient step is always
    paid.
  \item \emph{Wrong acceleration factor in Algorithm~3.} The factor
    $\beta_k^{\mathrm{pre}}$ had been implemented with Algorithm~1's
    $a_k=\nrm{g}^2$. That formula is the special case for $d=-\nabla\theta$.
    In Algorithm~3, $d=-B^{-1}g$, so the general form of their Section~2.3
    (the acceleration of Andrei~\cite{andrei2006acceleration}) is needed:
    \[
      \beta=\frac{-g^{\mathsf T}d}{(g_t-g)^{\mathsf T}d}\quad\text{if the denominator is positive, else }\beta=1,
    \]
    where $g_t$ is the gradient at the accepted trial point, as in
    \cref{alg:agd}.
    This form reduces exactly to Algorithm~1's when $d=-g$, so it is
    consistent with both statements in the paper. \emph{This defect caused
    the gap in outer iterations}: correcting it reduced P7 from 6 outer
    iterations to 4.
\end{enumerate}
The cumulative effect on P7 was 46 \EVDs\ $\to$ 40 (defects 1--2) $\to$ 43
(exact equation~(6)) $\to$ 32 (correct acceleration factor), against a
published 31. (\cref{tab:agd-full} shows only the GLL-guess, LF-guess, and
final columns; the 46-\EVD\ starting point and the 43-\EVD\ intermediate,
with the exact eq.~(6) but the uncorrected acceleration factor, are not
tabled.)

\paragraph{The remaining P7 residual.}
See \cref{sec:repro}. P7 matches in outer iterations (4) and residual norm
(21.0339) and differs by one backtracking decision (7 against 8). That decision
accounts for the whole difference in \EVDs, through the early-termination
saving. We attribute it to last-bit differences between the MATLAB and Julia
eigensolvers. P8 and P9 reproduce every count exactly on the same code path.
No further parameter search was carried out.

\paragraph{Conditions of use (all met in this paper).}
(i) The code and this record are published with the results. (ii) The P7
residual is disclosed. (iii) Published counts are printed beside local counts
wherever the comparison appears (\cref{tab:agd-verify,tab:agd-full}).

\paragraph{Checks against other published results.}
Beyond the AGD-SDAJ gate, every solver was run on invalid correlation matrices
with published solutions or work counts
(\cref{tab:external}). The matrices are those of the \texttt{CORRINV}
collection~\cite{higham2022anymatrix}, taken from its repository rather than
retyped. On the ten small matrices (\texttt{high02}, \texttt{tec03},
\texttt{bhwi01}, \texttt{mmb13}, \texttt{fing97}, \texttt{beyu11},
\texttt{tyda99r1}--\texttt{r3}, \texttt{usgs13}), run at tolerance $10^{-10}$,
the five solvers of the original comparison return valid correlation
matrices (Anderson-APM, added later, is checked in \cref{app:anderson}) whose
distances
$\fro{G-X}$ agree to at least eight decimal places and whose matrices agree
with Newton-SIN-BH's to within $2\times10^{-10}$, except unmodified AGD-SDAJ on
\texttt{mmb13} ($8.7\times10^{-10}$), which stops at its budget. The three
large matrices (\texttt{cor1399}, \texttt{cor3120} and \texttt{bccd16}) were
compared only at the published tolerances.

\begin{table}[h]
\centering
\caption{Work counts reproduced from the literature. Dykstra-APM uses the
stopping rule of the released code \texttt{nearcorr\_new.m}
of~\cite{higham2016anderson}, tolerance $n\,\mathrm{eps}$; run for the same
number of iterations, our implementation returns the same iterate, bit for
bit on nine matrices and to $1.4\times10^{-15}$ on \texttt{mmb13}. Newton-SIN-BH is compared with Tables 6.5 and 6.6
of~\cite{borsdorf2007newton}.}
\label{tab:external}
\small
\begin{tabular}{llll}
\toprule
solver & source & matrix & published / ours\\
\midrule
Dykstra-APM & \cite{higham2016anderson}, Table 1 & \texttt{tec03}, $n=4$ & 39 / 39 iterations\\
            & & \texttt{bhwi01}, $n=5$ & 27 / 27\\
            & & \texttt{mmb13}, $n=6$ & 801 / 801\\
            & & \texttt{fing97}, $n=7$ & 33 / 34$^{a}$\\
            & \cite{higham2016anderson}, Table 4 & \texttt{usgs13}, $n=94$ & 18 / 18\\
\midrule
Newton-SIN-BH & \cite{borsdorf2007newton} & \texttt{cor1399}, $10^{-7}n$ & 5 / 5 iterations\\
              & & \texttt{cor1399}, $n\,\mathrm{eps}$ & 7 / 7; $\fro{G-X}=21.033911$ / 21.033911\\
              & & \texttt{cor3120}, $10^{-7}n$ & 4 / 4\\
              & & \texttt{cor3120}, $n\,\mathrm{eps}$ & 6 / not reached$^{b}$; 5.443751 / 5.443751\\
\bottomrule
\end{tabular}
\end{table}
{\small $^{a}$At iteration 33 the stopping quantity is $1.12$ times the
tolerance $7\,\mathrm{eps}$, so the extra iteration is decided in the last
bits. $^{b}$The residual floor near $10^{-12}$ lies above
$n\,\mathrm{eps}=6.9\times10^{-13}$ with our eigensolver; the solution itself
matches to all published digits.\par}

On three of these matrices unmodified AGD-SDAJ shows the effect of
\cref{sec:trap2}: on \texttt{mmb13} it stalls (budget of 200\,000 \EVDs\
exhausted, against 210 for AGD-SDAJ-BH), and on \texttt{tyda99r2} and
\texttt{bhwi01} it converges, but slowly (549 against 35, and 66 against 21).

\paragraph{The Borsdorf--Higham adaptation.}
AGD-SDAJ-BH changes only the monotone Armijo test of Algorithm~1, line~5
(\cref{alg:agd}). When
$|\theta_t-\theta|<\gamma u(1+|\theta_t|+|\theta|)$ with $\gamma=100$ and
$u=2^{-53}$, a trial that reduces $\nrm{\nabla\theta}$ is accepted (weaker than
the Newton safeguard's $(1-\mubh)$ reduction by design: near the precision floor
a gradient step only inches the gradient down, so the stronger test never
fires). The
nonmonotone rule~\eqref{eq:hh6} is unchanged. This adaptation is ours.

\section{Artifact manifest}\label{app:artifact}

The generator, the solvers, the analysis scripts and every result file behind
every table are available at
\url{https://github.com/Realosunboy6/ncm-exact-benchmark}. Paths below are
relative to the artifact root. An archived snapshot carrying a permanent
identifier will be deposited on acceptance.

Instance binaries are not stored in full, since they reach 1\,GB at $n=500$;
the manifests record a hash of every instance used instead. The generator is
seeded, but it takes a basis for a degenerate eigenspace from LAPACK, so the
bits it produces depend on the LAPACK build: on another build a regenerated
set has the same parameters and passes the same screening, but is not
hash-identical to the one used here. The manifests and the screening logs are
therefore the authoritative record of the instances. A 12-instance subset at
$n=100$ (2\,MB) ships with the artifact for smoke tests.

\begin{table}[H]
\centering
\small
\begin{tabular}{p{0.40\linewidth}p{0.54\linewidth}}
\toprule
file & contents\\
\midrule
\multicolumn{2}{l}{\emph{Instances}}\\
\texttt{code/gen\_degeneracy\_family.py} & generator and screening (\cref{sec:family})\\
\texttt{code/materialize\_exact\_ystars.py} & exact $\ys=-\diag(W)$ for a suite\\
\texttt{data/canonical\_n100/} & 12 instances with $G$, $\Xs$, $\ys$, manifest and cases\\
\texttt{data/manifests/} & SHA-256 manifests and screening logs for every set\\
\texttt{code/export\_real\_suite.py} & literature matrices from Anymatrix \texttt{CORRINV}\\
\texttt{code/export\_thesis\_kkt.py}, \texttt{export\_thesis\_suite.py} & equity-derived sets (\cref{sec:thesis,sec:us2105})\\
\midrule
\multicolumn{2}{l}{\emph{Solvers}}\\
\texttt{code/bench\_sbb\_dual.jl} & \texttt{sbb\_dual}, \texttt{newton\_ncm\_globalized}, \texttt{agd\_sdaj\_ncm}, \texttt{dykstra\_apm}, trajectory recorder, study drivers\\
\texttt{code/anderson\_apm.jl} & Anderson-APM, transcribed from Higham and Strabi\'c's released code\\
\texttt{code/agd\_sdaj.jl} & standalone AGD-SDAJ used for verification (\cref{app:agd})\\
\texttt{code/validation/} & checks against Higham (2002), Higham--Strabi\'c (2016), Borsdorf (2007), Huynh--Hwang (2025), including the Anderson-APM record of \cref{app:anderson}; logs in \texttt{results/validation/}\\
\texttt{code/highrank\_queue.sh}, \texttt{anderson\_rest\_queue.sh} & drivers for the runs added with Anderson-APM\\
\texttt{code/check\_reference\_accuracy.jl} & two independent references per matrix (\cref{sec:real})\\
\midrule
\multicolumn{2}{l}{\emph{Results and analysis}}\\
\texttt{code/reproduce.sh} & regenerates the analysis tables from \texttt{results/*.csv.gz} and prints which file holds which table; the timing tables are regenerated separately by \texttt{analyze\_timing.py}\\
\texttt{results/ranking\_kkt270\_v2.csv.gz} & one row per \EVD, the five original solvers, 270 instances at $n=100$\\
\texttt{results/ranking\_n500.csv.gz} & the same at $n=500$\\
\texttt{results/ranking\_highrank\_n500.csv.gz} & $r\in\{75,100,150\}$, Newton-SIN-BH, AGD-SDAJ-BH and SBB-Dual (\cref{tab:highrank})\\
\texttt{results/ranking\_highrank\_proj.csv.gz} & the same instances, Dykstra-APM and Anderson-APM\\
\texttt{results/ranking\_n500\_dykstra1200.csv.gz} & the 90 censored Dykstra-APM runs at $n=500$ rerun at a cap of 1200 (\cref{sec:acct-results})\\
\texttt{results/ranking\_*anderson*.csv.gz} & Anderson-APM on every other test set, run after the original studies\\
\texttt{results/ranking\_thesis*.csv.gz} & equity sets at $n=550$ (\cref{tab:thesiskkt,tab:thesispert})\\
\texttt{results/ranking\_us2105\_*.csv.gz} & equity sets at $n=2105$ (\cref{tab:us2105,tab:us2105pert})\\
\texttt{results/ranking\_real.csv.gz} & four literature matrices (\cref{tab:real})\\
\texttt{results/timing\_*.csv} & single-threaded timings\\
\texttt{code/analyze\_ranking.py} & \cref{tab:fwd,tab:cells,tab:feas,tab:n500,tab:thesispert} and the accounting numbers of \cref{sec:acct-results}\\
\texttt{code/analyze\_by\_rank.py} & \cref{tab:highrank,tab:thesiskkt,tab:us2105}\\
\texttt{code/analyze\_per\_instance.py} & \cref{tab:real}\\
\texttt{code/analyze\_stopping\_rule.py} & the cross-dimension comparison of \cref{sec:stopping}\\
\texttt{code/analyze\_highrank.py}, \texttt{analyze\_us2105.py}, \texttt{analyze\_confound.py} & bootstrap intervals for \cref{tab:highrank}, \cref{tab:us2105pert}, and the bulk-level pilot\\
\texttt{code/check\_scoring\_column.py} & the rescaled-iterate sensitivity check (\cref{sec:acct-results})\\
\texttt{code/analyze\_timing.py} & \cref{tab:primitive,tab:perevd,tab:native}\\
\bottomrule
\end{tabular}
\end{table}

\paragraph{Data not redistributed.}
The Anymatrix \texttt{CORRINV} matrices carry their own terms and are obtained
from their own repository. The two equity panels are built from vendor price
data that may not be redistributed; the exporters, the selection rules and all
result files are included, the raw prices are not.

{\sloppy
\paragraph{Figure data.} The files \texttt{paper/figures/data/*.dat} were
extracted from the result files above. The file \texttt{evd\_vs\_eps.dat}
applies the cost-to-target rule of \texttt{analyze\_ranking.py}
(accepted-only, reach and hold) on a half-decade grid of $\varepsilon$. At the
five decade targets it reproduces \cref{tab:fwd}, up to that table's rounding
of half-integer medians.\par}


% refs.bib = paper2-lit.bib from the literature audit (26 Crossref-verified
% entries), merged 2026-09-11; the earlier draft entries were superseded.
\bibliographystyle{plain}
\bibliography{refs}

\end{document}
